\documentclass{article}

% Language setting
% Replace `english' with e.g. `spanish' to change the document language
\usepackage[english]{babel}

% Set page size and margins
% Replace `letterpaper' with `a4paper' for UK/EU standard size
\usepackage[letterpaper,top=2cm,bottom=2cm,left=3cm,right=3cm,marginparwidth=1.75cm]{geometry}

% Useful packages
\usepackage{amsmath, amssymb, array} % in the preamble
\usepackage{enumitem}
\usepackage{amsthm} % In your preamble
\usepackage{bussproofs}
\usepackage{amssymb}
\usepackage{graphicx}
\usepackage{tabularx} % Required for adjustable width tables
\usepackage{array}
\usepackage{tablefootnote}

\usepackage[colorlinks=true, allcolors=blue]{hyperref}

\title{SelVerIR: Self-Verifying Programming Language \\ Intermediate Representation Syntax and Semantics Document \\ v0.9}
\author{Yusuf Demir \\ Yağız Kaan Aydoğdu}

\begin{document}
\maketitle

% \begin{abstract}
% \end{abstract}

\newpage
\section{Syntax}

\[
\begin{array}{rcl}

\langle \mathit{IR} \rangle &::=& \{\langle \mathit{label} \rangle\ \langle \mathit{Command} \rangle\} \\[6pt]

\langle \mathit{Command} \rangle
&::=& \texttt{DECL} \; \langle basic\_type \rangle \; , \langle ident \rangle \;\mid\; \texttt{LDECL} \; \langle list\_type \rangle \; , \langle ident \rangle
\\[2pt]

&& \mid\; \texttt{PUSH} \; \langle imm \rangle
\\[2pt]

&& \mid\; \texttt{LOAD} \; \langle ident \rangle
\\[2pt]

&& \mid\; \texttt{STORE} \; \langle ident \rangle
\\[2pt]

&& \mid\; \texttt{LLOAD} \; \langle ident \rangle
\\[2pt]

&& \mid\; \texttt{LSTORE} \; \langle ident \rangle
\\[2pt]

&& \mid\; \texttt{ADD}
\\[2pt]

&& \mid\; \texttt{SUB}
\\[2pt]

&& \mid\; \texttt{MUL}
\\[2pt]

&& \mid\; \texttt{iDIV} \mid\; \texttt{fDIV}
\\[2pt]

&& \mid\; \texttt{LEN} \; \langle ident \rangle
\\[2pt]

&& \mid\; \texttt{EQ}
\\[2pt]

&& \mid\; \texttt{LT}
\\[2pt]

&& \mid\; \texttt{LE}
\\[2pt]

&& \mid\; \texttt{GT}
\\[2pt]

&& \mid\; \texttt{GE}
\\[2pt]

&& \mid\; \texttt{NEG}
\\[2pt]

&& \mid\; \texttt{AND}
\\[2pt]

&& \mid\; \texttt{OR}
\\[2pt]

&& \mid\; \texttt{XOR}
\\[2pt]

&& \mid\; \texttt{READ}
\\[2pt]

&& \mid\; \texttt{WRITE} \;\mid\; \texttt{LWRITE}
\\[2pt]

&& \mid\; \texttt{WRITELN} \;\mid\; \texttt{LWRITELN}
\\[2pt]

&& \mid\; \texttt{JZ} \; \langle int \rangle
\\[2pt]

&& \mid\; \texttt{GOTO} \; \langle int \rangle
\\[2pt]

&& \mid\; \texttt{CSCOPE} \;\mid\; \texttt{PSCOPE}
\\[2pt]

&& \mid\; \texttt{FUNCENV} \; \langle ident \rangle, \langle type \rangle
\\[2pt]

&& \mid\; \texttt{CALL} \; \langle ident \rangle
\\[2pt]

&& \mid\; \texttt{RET}
\\[2pt]

&& \mid\; \texttt{NOOP}
\\[15pt]

\langle \mathit{basic\_type} \rangle &::=& \texttt{INT} \;\mid\; \texttt{FLOAT}\\[6pt]

\langle \mathit{list\_type} \rangle &::=& \texttt{LIST[INT]} \;\mid\; \texttt{LIST[FLOAT]}  \\[6pt]

\langle \mathit{type} \rangle &::=& \langle \mathit{basic\_type} \rangle \;\mid\; \langle \mathit{list\_type} \rangle \\[6pt]

\langle \mathit{ident} \rangle &::=& \text{[A-Za-z\_][A-Za-z0-9\_]*} \\[6pt]

\langle \mathit{label} \rangle &::=& [0\texttt{-}9]^+ \texttt{:}\\[6pt]

\langle \mathit{int} \rangle &::=& [0\texttt{-}9]^+ \\[6pt]

\langle \mathit{float} \rangle &::=& [0\texttt{-}9]^+\texttt{.}[0\texttt{-}9]^+ \\[6pt]

\langle \mathit{imm} \rangle &::=& \langle \mathit{int} \rangle \;\mid\; \langle \mathit{float} \rangle \\[6pt]

\end{array}
\]

\newpage
\section{Semantics}

\subsection{Mathematical notions}
The \texttt{SelVeri} intermediate representation, \texttt{SelVerIR} is executed on a \textit{stack-based} abstract machine (\textbf{AM}) and aims to be provably correct implementation of the high-level \texttt{SelVeri} language.

\subsubsection{Stack}
At any moment of the execution of a \texttt{SelVerIR} program, its stack, denoted as $e$, is a an ordered finite sequence consisting elements are from \textbf{Val}, which will be used in the flow of the program in the \textit{LIFO} order. As expected, $e$ supports a variety of \texttt{PUSH} and (indirectly) \texttt{POP} operations. For convenience, the empty stack is denoted separately, as $\epsilon$. Furthermore, other instructions like \texttt{LOAD}, \texttt{LLOAD}, \texttt{STORE}, \texttt{LSTORE} can perform (possible multiple) push and pop operations on the stack.

\subsubsection{Program counter}
Intuitively, one can think of any \texttt{SelverIR} program, $P$ as an ordered list of commands. Program counter, $pc$, represents the index of the command that is to be executed next. Control-flow statements such as \texttt{GOTO} and \texttt{JZ} can alter the program counter in a non-sequential manner. Function calls using \texttt{CALL} instruction also make use of $pc$ internally.

\subsubsection{Return address}
A return address, $ra$, is the index of the command that is to be executed next, just after current function call ends. Instead of storing one $ra$ at any given time, the abstract machine manages a separate stack called $ras$, containing return addresses in LIFO order, to be able to support nested function calls.

\subsubsection{IR-configuration}
An IR-configuration is the representation of the execution state in the form of an 8-tuple 
\[
    \langle c,\, e,\, \sigma,\, s,\, pc,\, ras,\, \iota, \Omega \rangle
\] where
\begin{itemize}
    \item $c$ is the list of commands representing the remaining part of the program \texttt{SelverIR} program, i.e, $c = P[pc:]$.

    \item $e$ is the stack of the execution state,

    \item $\sigma$ is the program state,

    \item $s$ is the scope,

    \item $pc$ is the program counter.

    \item $ras$ is the return-address-stack.

    \item $\iota$ is the unread input stream.

    \item $\Omega$ is the output trace produced so far.
    
\end{itemize}

\newpage
\subsection{Small-step operational semantics of \texttt{SelverIR} commands}
Let $P$ be any \texttt{SelVerIR} program throughout the section \textbf{2.2}. For simplicity, we omit explicitly writing last 3 elements of the IR-configuration throughout the document, including the correctness proof at section \textbf{4}. Reader must be aware of the fact that these remain unchanged unless explicitly specified. Note that any configuration which is not explicitly stated in the tables above leads to the erroneous state, $\hat{\sigma}$ and increments $pc$ by exactly 1.

\subsubsection{Small-step operational semantics of state-related commands}

\begin{table}[htb!]
\setlength{\tabcolsep}{0pt} % Remove default padding to use full \textwidth
\centering
\renewcommand{\arraystretch}{2.2} 
\large
\begin{tabular*}{\textwidth}{@{\extracolsep{\fill}} >{\(}l<{\)} @{\quad \triangleright \quad} >{\(}l<{\)} @{}}
\hline\hline

\langle \texttt{DECL}\;t\, x : c ,\, e,\, \sigma,\, s,\, pc \rangle & 
\begin{cases}
\langle c,\, e,\, \sigma[x \mapsto \textbf{0}],\, s[x \mapsto t],\, pc+1 \rangle \\ \qquad \text{if } x \in \textbf{Var} \;\land\; t \in \{\texttt{INT}, \texttt{FLOAT}\} \;\land\; s(x) = \textbf{Void}\\
\end{cases} \\

\langle \texttt{LDECL}\;t\,x : c ,\, n : e,\, \sigma,\, s,\, pc \rangle & 
\begin{cases}
\langle c,\, e,\, \sigma[x \mapsto \textbf{0}],\, s[x \mapsto \texttt{LIST[INT, $n$]}],\, pc+1 \rangle \\ \qquad \text{if } x \in \textbf{Var} \;\land\; t = \texttt{LIST[INT]}]\;\land\; s(x) = \textbf{Void}\\

\langle c,\, e,\, \sigma[x \mapsto \textbf{0}],\, s[x \mapsto \texttt{LIST[FLOAT, $n$]}],\, pc+1 \rangle \\ \qquad \text{if } x \in \textbf{Var} \;\land\; t = \texttt{LIST[FLOAT]}]\;\land\; s(x) = \textbf{Void}\\

\langle c,\, e,\, \hat{\sigma},\, s,\, pc+1 \rangle \qquad \text{otherwise}}\\
\end{cases} \\

\langle \texttt{PUSH}\; x : c,\,  e,\, \sigma,\, s,\, pc \rangle & \begin{cases}
    \langle c,\, \mathcal{V}(x) : e,\, \sigma,\, s,\, pc+1 \rangle \qquad \text{if $x \in \textbf{Imm} \land x$ is of basic type }   \\
    \langle c,\, size(\textbf{List[$t,d,\vec{n}$]}) :  \mathcal{V}(x) : e,\, \sigma,\, s,\, pc+1 \rangle \qquad \text{if $x \in \textbf{Imm} \land $ type of $x$ is \textbf{List[$t,d,\vec{n}$]} }   \\ 
\end{cases}  \\


\langle \texttt{LOAD}\; x : c,\,  e,\, \sigma,\, s,\, pc \rangle & \begin{cases}
\langle c,\, \sigma(x,s) : e,\, \sigma,\, s,\, pc+1 \rangle & \\ \qquad \text{if } s(x) \in \{\texttt{INT}, \texttt{FLOAT}\} \land n=1 \land v_1 \in \mathbb{R} \land x \in \textbf{Var} \\[2pt]
\langle c,\, n : \sigma(x\texttt{[0]},s) : \dots : \sigma(x\texttt{[n-1]},s): e,\, \sigma,\, s,\, pc+1 \rangle & \\ \qquad \text{if } s(x) = \texttt{LIST[$t, n$]} \land x \in \textbf{Var} \\[2pt]
\langle c,\, e,\, \hat{\sigma},\, s,\, pc+1 \rangle \qquad \text{otherwise}
\end{cases} \\


\langle \texttt{STORE}\; x : c,\, v_0 : v_1 : \dots : v_n : e,\, \sigma,\, s,\, pc \rangle & \begin{cases}
\langle c,\, e,\, \sigma[x \mapsto v_0],\, s,\, pc+1 \rangle & \\ \qquad \text{if } s(x) \in \{\texttt{INT}, \texttt{FLOAT}\} \land n=0 \land v_0 \in \mathbb{R} \land x \in \textbf{Var} \\[2pt]
\langle c,\, e,\, \sigma[x \mapsto [v_1, \dots, v_n]],\, s,\, pc+1 \rangle & \\ \qquad \text{if } s(x) = \texttt{LIST[$t, v_0$]}\land v_1, \dots v_n \in \mathbb{R} \land x \in \textbf{Var} \\[2pt]
\langle c,\, e,\, \hat{\sigma},\, s,\, pc+1 \rangle \qquad \text{otherwise} \\
\end{cases}\\


\langle \texttt{LLOAD}\;x : c,\, i : e,\, \sigma,\, s,\, pc \rangle & \begin{cases}
\langle c,\, \sigma(x\texttt{[$i$]}, s) : e,\, \sigma,\, s,\, pc+1 \rangle & \\ \quad \text{if } s(x) = \texttt{LIST[$t,n$]} \land 0 \leq i \leq (n-1) & \\ \qquad  \land\; i,v \in \mathbb{N} \land x \in \textbf{Var} \\[2pt]
\langle c,\, e,\, \hat{\sigma},\, s,\, pc+1 \rangle \qquad \text{otherwise} \\
\end{cases}\\


\langle \texttt{LSTORE}\;x : c,\, i : v : e,\, \sigma,\, s,\, pc \rangle & \begin{cases}
\langle c,\, e,\, \sigma[x\texttt{[$i$]} \mapsto v],\, s,\, pc+1 \rangle & \\ \quad \text{if } s(x) = \texttt{LIST[$t,n$]} \land 0 \leq i \leq (n-1) & \\ \qquad  \land\; i,v \in \mathbb{N} \land x \in \textbf{Var} \\[2pt]
\langle c,\, e,\, \hat{\sigma},\, s,\, pc+1 \rangle \qquad \text{otherwise} \\
\end{cases}\\


\hline\hline
\end{tabular*}
\caption{Operational semantics of state-related commands}
\label{tab:am-sos-state}
\end{table}

\newpage
\subsubsection{Small-step operational semantics of arithmetic commands}

\begin{table}[htb!]
\setlength{\tabcolsep}{0pt}
\centering
\renewcommand{\arraystretch}{2.2} 
\large
\begin{tabular*}{\textwidth}{@{\extracolsep{\fill}} >{\(}l<{\)} @{\quad \triangleright \quad} >{\(}l<{\)} @{}}
\hline\hline


\langle \texttt{ADD} : c ,\, v_1 : v_2 : e,\, \sigma,\, s,\, pc \rangle & \langle c,\, (v_1 + v_2) : e,\, \sigma,\, s,\, pc+1 \rangle \quad \text{if } v_1, v_2 \in \mathbb{R} \\

\langle \texttt{SUB} : c ,\, v_1 : v_2 : e,\, \sigma,\, s,\, pc \rangle & \langle c,\, (v_1 - v_2) : e,\, \sigma,\, s,\, pc+1 \rangle \quad \text{if } v_1, v_2 \in \mathbb{R} \\

\langle \texttt{MUL} : c ,\, v_1 : v_2 : e,\, \sigma,\, s,\, pc \rangle & \langle c,\, (v_1 \cdot v_2) : e,\, \sigma,\, s,\, pc+1 \rangle \quad \text{if } v_1, v_2 \in \mathbb{R} \\

\langle \texttt{iDIV} : c ,\, z_1 : z_2 : e,\, \sigma,\, s,\, pc \rangle & \begin{cases} 
\langle c,\, \lfloor z_1 / z_2 \rfloor : e,\, \sigma,\, s,\, pc+1 \rangle & \text{if } z_1, z_2 \in \mathbb{Z} \land z_2 \neq 0 \\
\langle c,\, \bot : e,\, \hat{\sigma},\, s,\, pc+1 \rangle & \text{otherwise}
\end{cases} \\

\langle \texttt{fDIV} : c ,\, r_1 : r_2 : e,\, \sigma,\, s,\, pc \rangle & \begin{cases} 
\langle c,\, r_1 / r_2 : e,\, \sigma,\, s,\, pc+1 \rangle & \text{if } r_1, r_2 \in \mathbb{R} \land r_2 \neq 0 \\
\langle c,\, \bot : e,\, \hat{\sigma},\, s,\, pc+1 \rangle & \text{otherwise}
\end{cases} \\

\langle \texttt{LEN}\; x : c ,\, e,\, \sigma,\, s,\, pc \rangle & \begin{cases} 
\langle c,\, 0 : e,\, \sigma,\, s,\, pc+1 \rangle & \text{if } s(x) \in \{\texttt{INT}, \texttt{FLOAT}\} \\
\langle c,\, \mathcal{A}(\texttt{\_$lst$\_len\_1}, \sigma, s) : e,\, \sigma,\, s,\, pc+1 \rangle & \text{if } s(x) = \texttt{LIST[t, $n$]} \\
\langle c,\, \bot : e,\, \hat{\sigma},\, s,\, pc+1 \rangle & \text{otherwise}
\end{cases} \\


\hline\hline
\end{tabular*}
\caption{Operational semantics of arithmetic commands}
\label{tab:am-sos-arith}
\end{table}

\subsubsection{Small-step operational semantics of boolean commands}

\begin{table}[htb!]
\setlength{\tabcolsep}{0pt}
\centering
\renewcommand{\arraystretch}{2.2} 
\large
\begin{tabular*}{\textwidth}{@{\extracolsep{\fill}} >{\(}l<{\)} @{\quad \triangleright \quad} >{\(}l<{\)} @{}}
\hline\hline

\langle \texttt{EQ} : c ,\, v_1 : v_2 : e,\, \sigma,\, s,\, pc \rangle & \langle c,\, (v_1 = v_2) : e,\, \sigma,\, s,\, pc+1 \rangle \quad \text{if } v_1, v_2 \in \mathbb{R} \\

\langle \texttt{LE} : c ,\, v_1 : v_2 : e,\, \sigma,\, s,\, pc \rangle & \langle c,\, (v_1 \leq v_2) : e,\, \sigma,\, s,\, pc+1 \rangle \quad \text{if } v_1, v_2 \in \mathbb{R} \\

\langle \texttt{LT} : c ,\, v_1 : v_2 : e,\, \sigma,\, s,\, pc \rangle & \langle c,\, (v_1 < v_2) : e,\, \sigma,\, s,\, pc+1 \rangle \quad \text{if } v_1, v_2 \in \mathbb{R} \\

\langle \texttt{GE} : c ,\, v_1 : v_2 : e,\, \sigma,\, s,\, pc \rangle & \langle c,\, (v_1 \geq v_2) : e,\, \sigma,\, s,\, pc+1 \rangle \quad \text{if } v_1, v_2 \in \mathbb{R} \\

\langle \texttt{GT} : c ,\, v_1 : v_2 : e,\, \sigma,\, s,\, pc \rangle & \langle c,\, (v_1 > v_2) : e,\, \sigma,\, s,\, pc+1 \rangle \quad \text{if } v_1, v_2 \in \mathbb{R} \\

\langle \texttt{AND} : c,\, b_1 : b_2 : e,\, \sigma,\, s,\, pc \rangle & \begin{cases}
\langle c,\, 1 : e,\, \sigma,\, s,\, pc+1 \rangle & \text{if } b_1 = 1 \land b_2 = 1 \\
\langle c,\, 0 : e,\, \sigma,\, s,\, pc+1 \rangle & \text{if } b_1 = 0 \lor b_2 = 0
\end{cases} \\

\langle \texttt{OR} : c,\, b_1 : b_2 : e,\, \sigma,\, s,\, pc \rangle & \begin{cases}
\langle c,\, 1 : e,\, \sigma,\, s,\, pc+1 \rangle & \text{if } b_1 = 1 \lor b_2 = 1 \\
\langle c,\, 0 : e,\, \sigma,\, s,\, pc+1 \rangle & \text{if } b_1 = 0 \land b_2 = 0
\end{cases} \\

\langle \texttt{XOR} : c,\, b_1 : b_2 : e,\, \sigma,\, s,\, pc \rangle & \begin{cases}
\langle c,\, 1 : e,\, \sigma,\, s,\, pc+1 \rangle & \text{if } b_1 \neq b_2 \\
\langle c,\, 0 : e,\, \sigma,\, s,\, pc+1 \rangle & \text{if } b_1 = b_2
\end{cases} \\

\langle \texttt{NEG} : c,\, b : e,\, \sigma,\, s,\, pc \rangle & \begin{cases}
\langle c,\, 1 : e,\, \sigma,\, s,\, pc+1 \rangle & \text{if } b = 0 \\
\langle c,\, 0 : e,\, \sigma,\, s,\, pc+1 \rangle & \text{if } b = 1
\end{cases} \\
\hline\hline
\end{tabular*}
\caption{Operational semantics of boolean commands}
\label{tab:am-sos-bool}
\end{table}

\newpage
\subsubsection{Small-step operational semantics of I/O commands}

\begin{table}[htb!]
\setlength{\tabcolsep}{0pt}
\centering
\renewcommand{\arraystretch}{2.2} 
\large
\begin{tabular*}{\textwidth}{@{\extracolsep{\fill}} >{\(}l<{\)} @{\quad \triangleright \quad} >{\(}l<{\)} @{}}
\hline\hline

\langle \texttt{READ} : c ,\, e,\, \sigma,\, s,\, pc,\, ras,\, i : \iota,\, \Omega \rangle & 
{\renewcommand{\arraystretch}{1.0}
\begin{array}[t]{@{}l@{}}
\langle c,\, \mathcal{V}(i): e,\, \sigma,\, s,\, pc+1,\, ras,\, \iota,\, \Omega \rangle \\
\text{if $i \in \textbf{Imm} \cap \{\textbf{Int} \cup \textbf{Float}\}$}
\end{array}} \\

\langle \texttt{WRITE} : c ,\, x : e,\, \sigma,\, s,\, pc,\, ras,\, \iota,\, \Omega \rangle & 
\langle c,\, e,\, \sigma,\, s,\, pc+1,\, ras,\, \iota,\, \mathcal{A}(x, \sigma, s) : \Omega \rangle \\

\langle \texttt{WRITELN} : c ,\, x : e,\, \sigma,\, s,\, pc,\, ras,\, \iota,\, \Omega \rangle & 
\langle c,\, e,\, \sigma,\, s,\, pc+1,\, ras,\, \iota,\, \texttt{nl} : \mathcal{A}(x, \sigma, s) : \Omega \rangle \\
\hline\hline
\end{tabular*}
\caption{Operational semantics of I/O commands}
\label{tab:am-sos-ctrl}
\end{table}

\texttt{nl} denotes the newline character in the usual sense and is handled by the abstract machine and its interaction with the operating system. It is included for convenience and SelVerIR still does not have any direct interaction with the IO stream or char types. Write mode does not affect the operational semantics in any way.

\subsubsection{Small-step operational semantics of control-flow commands}

\begin{table}[htb!]
\setlength{\tabcolsep}{0pt}
\centering
\renewcommand{\arraystretch}{2.2} 
\large
\begin{tabular*}{\textwidth}{@{\extracolsep{\fill}} >{\(}l<{\)} @{\quad \triangleright \quad} >{\(}l<{\)} @{}}
\hline\hline
\langle \texttt{JZ}\; z: c,\, b : e,\, \sigma,\, s,\, pc \rangle & \begin{cases}
\langle c,\, e,\, \sigma,\, s,\, pc+1 \rangle & \text{if } b = 1 \\
\langle P[\mathcal{I}(z):],\, e,\, \sigma,\, s,\, \mathcal{I}(z) \rangle & \text{if } b = 0 \; \land \; z \in \textbf{Imm $\cap$ Int}
\end{cases} \\

\langle \texttt{GOTO}\; z : c ,\, e,\, \sigma,\, s,\, pc \rangle & \langle P[\mathcal{I}(z):],\, e,\, \sigma,\, s,\, \mathcal{I}(z) \rangle \quad \text{if } z \in \textbf{Imm $\cap$ Int} \\

\langle \texttt{CALL}\; f : c ,\, e,\, \sigma,\, s,\, pc,\, ras \rangle & \langle P[pc_f:],\, e,\, \sigma,\, s,\, pc_f,\, (pc+1) : ras \rangle \quad \text{if } f \in \textbf{Func} \land \Gamma(f) \neq \emptyset \\

\langle \texttt{CSCOPE} : c ,\, e,\, \sigma,\, s,\, pc \rangle & \langle c,\, e,\, \tilde{\sigma_c},\, \tilde{s_c},\, pc+1 \rangle \quad \text{where $\sigma$ and $s$ are the parents of $\sigma_c$ and $s_c$}  \\

\langle \texttt{PSCOPE} : c ,\, e,\, \sigma_c,\, s_c,\, pc \rangle & \langle c,\, e,\, \sigma,\, s,\, pc+1 \rangle \quad \text{where $\sigma$ and $s$ are the parents of $\sigma_c$ and $s_c$}  \\

\langle \texttt{FUNCENV}\; f, t_r : c ,\, e,\, \sigma,\, s,\, pc \rangle & \langle c,\, e,\, \tilde{\sigma}',\, \tilde{s}',\, pc+1 \rangle \quad \text{where $f$ names the current function and $\sigma$ and $s$ is of the caller}  \\

\langle \texttt{RET} : c ,\, v : e,\, \sigma',\, s',\, pc ,\, ra : ras \rangle & \langle P[ra:],\, e,\, \sigma[\texttt{retvar} \mapsto v],\, s[\texttt{retvar} \mapsto s'(v)],\, ra ,\, ras \rangle \quad \text{if the return-type of the function called is equal to $s'(v)$ and $\sigma$ and $s$ is of the caller} \\

\langle \texttt{NOOP} : c ,\, e,\, \sigma,\, s,\, pc \rangle & \langle c,\, e,\, \sigma,\, s,\, pc+1 \rangle \\

\langle i : c ,\, e,\, \sigma,\, s,\, pc \rangle & \langle c,\, e,\, \hat{\sigma},\, s,\, pc+1 \rangle \quad \text{where $i$ is any command} \\
\hline\hline
\end{tabular*}
\caption{Operational semantics of control-flow commands}
\label{tab:am-sos-ctrl}
\end{table}

As one can notice from the tables above, we do not have type-specific commands for \texttt{ADD}, \texttt{SUB} and \texttt{MUL} or boolean operators as these operations are closed under $\mathbb{Z}$ \footnote{Although syntactically different in \texttt{SelVeri} high level language; semantically, we do not interpret any integer differently than its real number counterpart.}. However, since this is not the case for division, we have separate instructions for integer and float division, namely \texttt{iDIV} and \texttt{fDIV}.

\newpage
\section{Determinism of \texttt{SelVerIR}}

Before proceeding to the compilation section, we need to guarantee that the \texttt{SelVerIR} abstract machine serves as a mathematically sound target for the compiler. To achieve this we must prove that its execution is strictly deterministic. This requires showing that the small-step operational semantics transition relation ($\triangleright$) for the intermediate representation is a partial function. \\

\textbf{Theorem (Determinism of SelVerIR):} For any IR-configuration $\gamma = \langle c,\, e,\, \sigma,\, s,\, pc,\, ras,\, \iota,\, \Omega \rangle$, if $\gamma \triangleright \gamma_1$ and $\gamma \triangleright \gamma_2$, then $\gamma_1 = \gamma_2$. \\

\textbf{Proof:} The proof proceeds by case analysis on the current command to be executed, which is the head of the command list $c$. We must demonstrate that for any configuration, exactly one transition rule from the operational semantics tables applies. \\

\textit{Step 1: Uniqueness of Command Matching} \\
The execution of the abstract machine is strictly driven by the command at $c[0]$. Because the syntax of \texttt{SelVerIR} commands (e.g., \texttt{ADD}, \texttt{LOAD}, \texttt{JZ}) is unambiguous, the configuration will match the left-hand side of at most one command's semantic definition. \\

\textit{Step 2: Mutual Exclusivity of Command Premises} \\
For commands that define multiple possible resulting states, we must show that their branching conditions are mutually exclusive:
\begin{itemize}
    \item \textbf{State-related Commands (\texttt{DECL}, \texttt{LOAD}, \texttt{STORE}):} These commands branch strictly based on the static type of the identifier in the scope ($s(x)$) or the length of the list being accessed. Since a variable $x$ cannot simultaneously map to multiple disparate types or lengths in $s$, these conditions are mutually exclusive.
    \item \textbf{Arithmetic and Boolean Commands:} Instructions like \texttt{iDIV} and \texttt{fDIV} branch on whether the divisor is zero or non-zero ($z_2 = 0$ vs. $z_2 \neq 0$). Boolean commands like \texttt{AND}, \texttt{OR}, and \texttt{NEG} branch on explicit, mutually exclusive truth values ($1$ or $0$) popped from the stack.
    \item \textbf{Control-flow (\texttt{JZ}):} The \texttt{JZ} instruction branches on whether the condition value $b$ popped from the stack is exactly $1$ or exactly $0$. It is mathematically impossible for both to be true simultaneously.
\end{itemize}

\textit{Step 3: Deterministic Error Fallback} \\
For any valid command where the current stack $e$ or scope $s$ fails to satisfy the specific preconditions listed in the semantic definitions (e.g., a type mismatch, stack underflow, or an invalid list index), the semantics explicitly state that the configuration deterministically transitions to the erroneous state $\hat{\sigma}$ and increments the program counter $pc$ by exactly 1. \\

Because every valid \texttt{SelVerIR} command maps to exactly one deterministic outcome based on the state of the stack, scope, and program variables, and because any undefined configuration maps strictly to a single defined error state, at most one rule can apply to any IR-configuration. With a simple induction, this rule can be generalized to any \texttt{SelVerIR} program $c$ of finite length. Therefore, $\gamma_1 = \gamma_2$, proving that the \texttt{SelVerIR} is deterministic. $\blacksquare$

\newpage
\section{Compilation}

Let \textbf{Code} and \textbf{CodeIR} be the sets of all \texttt{SelVeri} and \texttt{SelVerIR} programs, respectively. Using these sets, we will inductively define a compiler function:

\[
\mathcal{C}: \textbf{Code} \rightarrow \textbf{CodeIR}
\]

\subsection{Compilation of lists}

Unlike high-level \texttt{SelVeri} language, \texttt{SelVerIR} does not support multi-dimensional lists, too simplify the maintance of stack during program execution. However, as the lists are interpreted as a ‘batch of variables’, we may flatten any multi-dimensional list into a one-dimensional list of its basic type. We call this technique \textit{list flattening} and to mathematically formalize it, we define the following function:


\[
\mathit{LF}: \textbf{Var} \times \textbf{AExp}^m \rightarrow \textbf{AExp}
\] where $m$ is the dimension of the input list variable.

\[
\mathit{LF} (lst\;; i_1, \dots, i_d) 
= \footnote{Because \mathit{LF} constructs an arithmetic expression rather than computing a raw integer, the $\Sigma$ and $\Pi$ symbols here denote the syntactic generation of chained $\texttt{+}$ and \texttt{$\ast$} expressions}
\sum_{j=1}^{d} i_j \left( \prod_{k=j+1}^{d} \texttt{\_$lst$\_len\_$k$} \right)
\] 

$\mathit{LF}$ is obviously well-defined, however we need to prove that $\mathit{LF}$ is injective on indices $(i_j)$ as well to ensure that the compilation is correct.

\begin{proof}
Even though not explicitly stated in its formal definition, all $i_j$ represent indices and $n_k$ represent length of the dimension $k$. As $\mathit{LF}$ will solely be used to flatten a valid indexing of an $d$-dimensional list, we assume standard 0-based indexing, meaning each index $i_k$ is bounded by $0 \le i_k < n_k$.


Let $W_j = \prod_{k=j+1}^{d} n_k$ represent the weight of the $j$-th dimension. Note that $W_m = 1$, and $W_{j-1} = n_j W_j$.

We can rewrite the index calculation as:
$$\mathit{LF}(lst\;; i_1, \dots, i_d)= i_1 W_1 + i_2 W_2 + \dots + i_d W_d$$

To prove injectivity, we must show that if $\mathit{LF}(lst\;; i_1, \dots, i_d) = \mathit{LF}(lst\;; i'_1, \dots, i'_d)$, then $i_j = i'_j$ for all $1 \le j \le d$.

We can prove this by examining the maximum possible value of any suffix sum of the index calculation. The maximum value for the remaining dimensions from $r+1$ to $d$ occurs when every index is at its maximum possible value ($n_j - 1$):
$$\text{Max Sum} = \sum_{j=r+1}^{d} (n_j - 1) W_j$$

Because $W_{j-1} = n_j W_j$, this sum elegantly telescopes:
$$\text{Max Sum} = ((n_{r+1} - 1) W_{r+1}) + \dots + ((n_d - 1) W_d)$$
$$\text{Max Sum} = W_r - W_d = W_r - 1$$

This demonstrates a crucial property: the sum of all lower-dimensional components ($i_{r+1} W_{r+1} + \dots + i_d W_d$) is strictly less than the weight of the current dimension $W_r$.

Because the remainder is strictly smaller than the step size of the current dimension, it is impossible for a change in a lower dimension to overflow and alter a higher dimension. Therefore, two different sets of indices must produce different flat indices, proving that $\mathit{LF}$ is injective for $d$-dimensional lists. \\
\end{proof}


However, $\mathit{LF}$ only formalizes list-accesses in \texttt{SelVerIR}, but we must also formalize list declarations in a way compatible with list flattening. Let us define a helper function :

\[
\texttt{size\_IR} : \textbf{ListType} \rightarrow \textbf{AExp} \\

\texttt{size\_IR}($t$) = \begin{cases}
    \texttt{1} & \text{if $t \in \textbf{BasicType}$ } \\
    \Pi_{i=1}^d{a_i} & \text{if $t = \textbf{List[$t', d, \vec{a}$]}$ for some $t' \in \textbf{BasicType}$ }
\end{cases}
\] Similar to $\mathit{LF}$, \texttt{size\_IR} outputs a syntctic arithmetic expression and the use of $\Pi$ is a shorthan for represeting a chaing of $d$-many $\ast$ operations. \\

Observe that the valuation of \texttt{size\_IR}($t$) is equal $size(t)$. With the help of $\mathit{LF}$ and \texttt{size\_IR}, we will be able to correctly define the compilation of list-accesses and list-declarations in \textbf{sections 4.2} and \textbf{4.4}.


\newpage
\subsection{Compilation of arithmetic expressions}

In order to formally define the compilation of arithmetic expressions, we define a helper function:

\[
\mathcal{CA}: \textbf{AExp} \rightarrow \textbf{CodeIR}
\]

\begin{table}[h!]
\centering
\renewcommand{\arraystretch}{1.5} 

% \resizebox dynamically scales the box to exactly the width of the line.
\resizebox{\linewidth}{!}{%
\fbox{%
$\begin{array}{ll}
\mathcal{CA}(x) = \texttt{PUSH}\;x \qquad \text{if } x \in \textbf{Imm}\\[1ex]

\mathcal{CA}(x) = \texttt{LOAD}\; x \qquad \text{if } x \in \textbf{Var} \\[1ex]

\mathcal{CA}(lst\texttt{[$i_1$]}\dots\texttt{[$i_d$]}) = \mathcal{CA}(\mathit{LF}(lst; i_1, \dots, i_d)) : \texttt{LLOAD}\;lst \\ \qquad \text{if $lst$ is an $d$-dimensional list} \\[1ex]

\mathcal{CA}(\texttt{len($x$)}) = \texttt{LEN}\; x \\[1ex]

\mathcal{CA}(\texttt{read()}) = \texttt{READ} \\[1ex]

\mathcal{CA}(a_1 + a_2) = \mathcal{CA}(a_2): \mathcal{CA}(a_1) : \texttt{ADD} & \\[1ex]

\mathcal{CA}(a_1 \ast a_2) = \mathcal{CA}(a_2): \mathcal{CA}(a_1) : \texttt{MUL} & \\[1ex]

\mathcal{CA}(a_1 - a_2) = \mathcal{CA}(a_2): \mathcal{CA}(a_1) : \texttt{SUB} & \\[1ex]

\mathcal{CA}(a_1 / a_2) = \begin{cases}
\mathcal{CA}(a_2): \mathcal{CA}(a_1) : \texttt{iDIV} & \parbox[t]{0.35\textwidth}{\raggedright  if $(a_1 / a_2)$ has no subterm of type \footnotemark \textbf{Float}} \\[3ex]
\mathcal{CA}(a_2): \mathcal{CA}(a_1) : \texttt{fDIV} & \parbox[t]{0.35\textwidth}
{\raggedright  if $(a_1 / a_2)$ has at least one subterm of type \textbf{Float}} \\[1ex]
\end{cases}


\end{array}$%
}%
} % <-- closes \resizebox
\end{table}
\footnotetext{The types of subterms are determined by the compiler during the compile time.}


\newpage
\subsection{Compilation of boolean expressions}

In order to formally define the compilation of boolean expressions, we define a helper function:

\[
\mathcal{CB}: \textbf{BExp} \rightarrow \textbf{CodeIR}
\]

\begin{table}[h!]
\centering
\renewcommand{\arraystretch}{1.8} 

\fbox{%
$\begin{array}{ll}
\mathcal{CB}(\texttt{true}) = \texttt{PUSH 1} \\[1ex]

\mathcal{CB}(\texttt{false}) = \texttt{PUSH 0} \\[1ex]

\mathcal{CB}(a) = \mathcal{CA}(a) : \texttt{PUSH 0} : \texttt{EQ} : \texttt{NEG} 
\qquad \text{if } a \in \textbf{AExp}
\\[1ex]

\mathcal{CB}(a_1 \;\texttt{=}\; a_2) = \mathcal{CA}(a_2) : \mathcal{CA}(a_1) : \texttt{EQ} & \\[1ex]

\mathcal{CB}(a_1 \;\texttt{<=}\; a_2) = \mathcal{CA}(a_2) : \mathcal{CA}(a_1) : \texttt{LE} & \\[1ex]

\mathcal{CB}(\texttt{not}\; b) = \mathcal{CB}(b) : \texttt{NEG} & \\[1ex]

\mathcal{CB}(b_1 \;\texttt{and}\; b_2) = \mathcal{CB}(b_2) : \mathcal{CB}(b_1) : \texttt{AND} & \\[1ex]
\end{array}$%
}
\end{table}


\newpage
\subsection{Compilation of statements}

Finally, with the help of $\mathcal{CA}$ and $\mathcal{CB}$, we are ready to properly define the compilation function:

\[
\mathcal{C}: \textbf{Code} \rightarrow \textbf{CodeIR}
\]

\begin{table}[h!]
\centering
\renewcommand{\arraystretch}{1.8} 

% \resizebox dynamically scales the box to exactly the width of the line.
\resizebox{\linewidth}{!}{%
\fbox{%
$\begin{array}{ll}

\mathcal{C}(x \;\texttt{:}\; t) = \texttt{DECL}\; t\;x \qquad \text{if } x \in \textbf{Var} \text{ and $t \in \textbf{BasicType}$}
\\[1.5ex]

\mathcal{C}(lst \;\texttt{:}\; \textbf{List[$t, d, \vec{a}$]})
\\ \quad =
\texttt{DECL INT}\;\texttt{\_$lst$\_len\_1} : \mathcal{CA}(a_1) : \texttt{STORE}\;\texttt{\_$lst$\_len\_1} : \dots : \texttt{DECL INT}\;\texttt{\_$lst$\_len\_$d$} : \mathcal{CA}(a_d) : \texttt{STORE}\;\texttt{\_$lst$\_len\_$d$} 
\\ \quad
: \mathcal{CA}(\texttt{size\_IR(} \textbf{List[$t,d,\vec{a}$]} \texttt{)} ) : \texttt{LDECL LIST[$t$]}\;lst
\\[1.5ex]

\mathcal{C}(x \;\texttt{:=}\; a) = \mathcal{CA}(a) : \texttt{STORE}\;x \qquad \text{if } x \in \textbf{Var} \text{ and $x$ is of a basic type}
\\[1.5ex]

\mathcal{C}(lst\texttt{[$i_1$]}\dots\texttt{[$i_d$]} \;\texttt{:=}\; a) = \mathcal{CA}(a) : \mathcal{CA}(\mathit{LF}(lst; i_1, \dots, i_d)) :  \texttt{LSTORE}\;lst \\ \qquad \text{if } lst \in \textbf{Var} \text{ and $lst$ is an $d$-dimensional list}
\\[1.5ex]

\mathcal{C}(\texttt{pass}) = \texttt{NOOP}
\\[1.5ex]

\mathcal{C}(S_1;S_2) = \mathcal{C}(S_1) : \mathcal{C}(S_2)
\\[1.5ex]

\mathcal{C}(\texttt{if } b \texttt{ then } S \texttt{ fi}) = 
\mathcal{CB}(b) : \texttt{JZ}\; pc_{\texttt{NOOP}} : \texttt{CSCOPE} : \mathcal{C}(S) : \texttt{PSCOPE} :\texttt{NOOP}
\\[1.5ex]

\mathcal{C}(\texttt{if } b \texttt{ then } S_1 \texttt{ else } S_2 \texttt{ fi}) 
\\ \quad = 
\mathcal{CB}(b) : \texttt{JZ}\; pc_{\texttt{CSCOPE}_2} : \texttt{CSCOPE} : \mathcal{C}(S_1) : \texttt{PSCOPE} : \texttt{GOTO}\; pc_{\texttt{NOOP}} : \texttt{CSCOPE} : \mathcal{C}(S_2) : \texttt{PSCOPE} : \texttt{NOOP}
\\[1.5ex]

\mathcal{C}(\texttt{while } b \texttt{ do } S \texttt{ od}) = 
\mathcal{CB}(b) : \texttt{JZ}\; pc_{\texttt{NOOP}} : \texttt{CSCOPE} : \mathcal{C}(S) : \texttt{PSCOPE} : \texttt{GOTO}\; pc_b : \texttt{NOOP}
\\[1.5ex]

\mathcal{C}(\texttt{function}\; f(x_1 : t_1 \dots, x_n : t_n) \,\texttt{->}\, t_r \,\texttt{::}\, S_f) 
\\ \quad = \texttt{FUNCENV}\; f, t_r : [\texttt{L}]\texttt{DECL}\; t_1\;x_1 : \texttt{STORE}\; x_1 :  \dots : [\texttt{L}]\texttt{DECL}\; t_n\;x_n : \texttt{STORE}\; x_n : \mathcal{C}(S_f)
\\[1.5ex]

\mathcal{C}(\texttt{call}\; f(a_1 \dots, a_n)) = 
\mathcal{CA}(a_n) : \dots : \mathcal{CA}(a_1) : \texttt{CALL}\; f
\\[1.5ex]


\mathcal{C}(\texttt{return}\; a) = 
\mathcal{CA}(a) : \texttt{RET}
\\[1.5ex]

\mathcal{C}(\texttt{write}(a)) = \left\{
    \begin{aligned}
        \mathcal{CA}(a) : \texttt{WRITE} \quad &\text{if } \mathcal{A}(a, \sigma, s) \in \textbf{BasicType} \\
        \mathcal{CA}(a) : \texttt{LWRITE} \quad &\text{if } \mathcal{A}(a, \sigma, s) \in \textbf{ListType}
    \end{aligned}
    \right.
\\[1.5ex]

\mathcal{C}(\texttt{writeline}(a)) = \left\{
    \begin{aligned}
        \mathcal{CA}(a) : \texttt{WRITELN} \quad &\text{if } \mathcal{A}(a, \sigma, s) \in \textbf{BasicType} \\
        \mathcal{CA}(a) : \texttt{LWRITELN} \quad &\text{if } \mathcal{A}(a, \sigma, s) \in \textbf{ListType}
    \end{aligned}
    \right.
\\[1.5ex]
\end{array}$%
}%
} % <-- closes \resizebox
\end{table}

Please note that, during the compilation phase, compiler labels each compiled command with a corresponding program counter, using a simple incrementing counter. Therefore, $pc$ variables seen in the compilation of the control-flow statements are actually replaced with syntactic integer-immediates.

Regarding the in-brackets parts of function definition compilation, a function parameter can be either a list (declared using \texttt{LDECL})  or a basic type (declared using \texttt{DECL}) To cover this up, our operational semantics ensure that \texttt{STORE} also pushes the size of a list to the stack and \texttt{LDECL} expects the size of the declared list to be on top of the stack. However, since in both cases (with or without in-brackets parts) the resulting program state and scope tuple is trivially equivalent, we omit this distinction between \texttt{DECL} and \texttt{LDECL} in the correctness proof in the next section.

Additionally, as high-level \texttt{SelVeri} languagae ensures that every possible execution of $S_f$ ends with a \texttt{return}-statement, one can certainly say that $\mathcal{C}(S_f)$ ends with a \texttt{RET} command. This fact will be important for our correctess proof and for convenience, it was named as the \textit{return rule}.

The \texttt{LWRITE} and \texttt{LWRITELN} instructions are needed due to complete seperation of compiler and interpreter modules. Interpreter has no way to check if the value on top of the stack represents a literal value, or the size of the list packed placed in the stack.

\section{Provably Correct Implementation}
Now that the compilation function is explicitly defined, we need to prove the correctness of our implementation; by showing that given any  \texttt{SelVeri} program $P$, we have
\[
\langle P,\, \tilde{\sigma},\, \tilde{s} \rangle \; \rightarrow^{\ast} \; \underline{\sigma}_f \quad \implies \quad 
\langle \mathcal{C}(P),\, \epsilon,\, \tilde{\sigma},\, \tilde{s},\, 0 \rangle \; \triangleright^{\ast} \; \underline{\sigma}_f
\]

\subsection{Correctness of expressions}

\subsubsection{Arithmetic}
Given any program stack $e$, non-erroneous program state $\sigma$, scope $s$, program counter $pc$, we have:
\[
\forall a \in \textbf{AExp} \,\cdot\, \langle \mathcal{CA}(a) : c,\, e,\, \sigma,\, s,\, pc \rangle \; \triangleright^{\ast} \;
\langle c,\, \mathcal{A}(a, \sigma, s) : e,\, \underline{\sigma},\, s,\, pc + l \rangle
\] where $l$ is the number of commands in $\mathcal{CA}(a)$.

\begin{proof}
We will prove this by structural induction on arithmetic expressions.

\begin{itemize}
    \item \textbf{case 1: } $a = v$ for some $v \in \textbf{Imm} \cap (\textbf{Int} \cup \textbf{Float})$

    \[
      \begin{aligned}
        \langle  \mathcal{CA}(a) : c,\, e,\, \sigma,\, s,\, pc \rangle 
        &= \langle \mathcal{CA}(v) : c,\, e,\, \sigma,\, s,\, pc \rangle \\
        &= \langle \texttt{PUSH}\;v   : c,\, e,\, \sigma,\, s,\, pc \rangle 
        && \text{(definition of $\mathcal{CA}$)} \\
        &\triangleright \langle c,\, \mathcal{V}(v) : e,\, \sigma,\, s,\, pc+1 \rangle 
        && \text{(operational semantics)} \\
        &= \langle c,\, \mathcal{A}(v, \sigma, s) : e,\, \sigma,\, s,\, pc+1 \rangle 
        && \text{(definition of $\mathcal{A}$)} \\
        &= \langle c,\, \mathcal{A}(a, \sigma, s) : e,\, \sigma,\, s,\, pc+1 \rangle \\
      \end{aligned}
  \]

  \item \textbf{case 2: } $a = v$ for some $v \in \textbf{Imm} \cap (\bigcup \textbf{ListType})$, i.e. $v$ is a list-immediate of type \textbf{List[$t, d, \vec{n}$]}

    \[
      \begin{aligned}
        \langle  \mathcal{CA}(a) : c,\, e,\, \sigma,\, s,\, pc \rangle 
        &= \langle \mathcal{CA}(v) : c,\, e,\, \sigma,\, s,\, pc \rangle \\
        &= \langle \texttt{PUSH}\;v   : c,\, e,\, \sigma,\, s,\, pc \rangle 
        && \text{(definition of $\mathcal{CA}$)} \\
        &\triangleright \langle c,\, size(\textbf{List[$t, d, \vec{n}$]}) : \mathcal{V}(v) : e,\, \sigma,\, s,\, pc+1 \rangle 
        && \text{(operational semantics)} \\
        &= \langle c,\, \mathcal{A}(v, \sigma, s) : e,\, \sigma,\, s,\, pc+1 \rangle 
        && \text{(definition of $\mathcal{A}$)} \\
        &= \langle c,\, \mathcal{A}(a, \sigma, s) : e,\, \sigma,\, s,\, pc+1 \rangle \\
      \end{aligned}
  \]

    \item \textbf{case 3: } $a = x$ for some $x \in \textbf{Var}$ and $x$ is of a basic type

    \[
      \begin{aligned}
        \langle  \mathcal{CA}(a) : c,\, e,\, \sigma,\, s,\, pc \rangle 
        &= \langle \mathcal{CA}(x) : c,\, e,\, \sigma,\, s,\, pc \rangle \\
        &= \langle \texttt{LOAD}\;x : c,\, e,\, \sigma,\, s,\, pc \rangle 
        && \text{(definition of $\mathcal{CA}$)} \\
        &\triangleright \langle c,\, \sigma(x, s) : e,\, \sigma,\, s,\, pc+1 \rangle 
        && \text{(operational semantics)} \\
        &= \langle c,\, \mathcal{A}(x, \sigma, s) : e,\, \sigma,\, s,\, pc+1 \rangle 
        && \text{(definition of $\mathcal{A}$)} \\
        &= \langle c,\, \mathcal{A}(a, \sigma, s) : e,\, \sigma,\, s,\, pc+1 \rangle \\
      \end{aligned}
  \]

  \item \textbf{case 4: } $a = lst$ for some $lst \in \textbf{Var}$ and the type of $lst$ is \textbf{List[$t, d, \vec{n}$]}

    Observe that the size of $lst$ is compiled into $size(\textbf{List[$t, d, \vec{n}$]})$

  \[
      \begin{aligned}
        \langle  \mathcal{CA}(a) : c,\, e,\, \sigma,\, s,\, pc \rangle 
        &= \langle \mathcal{CA}(lst) : c,\, e,\, \sigma,\, s,\, pc \rangle \\
        &= \langle \texttt{LOAD}\;lst : c,\, e,\, \sigma,\, s,\, pc \rangle 
        && \text{(definition of $\mathcal{CA}$)} \\
        &\triangleright \langle c,\, size(\textbf{List[$t, d, \vec{n}$]}) : \sigma(x, s) : e,\, \sigma,\, s,\, pc+1 \rangle 
        && \text{(operational semantics)} \\
        &= \langle c,\, \mathcal{A}(lst, \sigma, s) : e,\, \sigma,\, s,\, pc+1 \rangle 
        && \text{(definition of $\mathcal{A}$)} \\
        &= \langle c,\, \mathcal{A}(a, \sigma, s) : e,\, \sigma,\, s,\, pc+1 \rangle \\
      \end{aligned}
  \]

  \item \textbf{case 5: } $a = lst\texttt{[$i_1$]}\dots\texttt{[$i_d$]}$ for some $d$-dimensional $lst$

  Suppose $\mathit{LF}(lst; i_1, \dots, i_d)) = j$ for some $j \in \textbf{AExp}$ and assume for an induction on the correctness of the compilation of $j$. 

    \[
      \begin{aligned}
        \langle  \mathcal{CA}(a) : c,\, e,\, \sigma,\, s,\, pc \rangle 
        &= \langle \mathcal{CA}(lst\texttt{[$i_1$]}\dots\texttt{[$i_d$]}) : c,\, e,\, \sigma,\, s,\, pc \rangle \\
        &= \langle \mathcal{CA}(\mathit{LF}(lst; i_1, \dots, i_d)) : \texttt{LLOAD}\;lst : c,\, e,\, \sigma,\, s,\, pc \rangle 
        && \text{(definition of $\mathcal{CA}$)} \\
        &= \langle \mathcal{CA}(j) : \texttt{LLOAD}\;lst : c,\, e,\, \sigma,\, s,\, pc \rangle 
        && \text{(assumption)} \\
        &\triangleright^\ast \langle \texttt{LLOAD}\;lst : c,\, \mathcal{A}(j, \sigma, s):  e,\, \sigma,\, s,\, pc + l_j \rangle 
        && \text{(induction hypothesis)} \\
        &\triangleright \langle c,\, \sigma(lst\texttt{[$\mathcal{A}(j, \sigma, s)$]}, s):  e,\, \sigma,\, s,\, pc + l_j + 1 \rangle
        && \text{(operational semantics)} \\
        &= \langle c,\, \sigma(lst\texttt{[$\mathcal{A}(\mathit{LF}(lst\;; i_1, \dots, i_d), \sigma, s)$]}, s):  e,\, \sigma,\, s,\, pc + l_j + 1 \rangle
        && \text{(definition of $\mathit{LF}$)} \\
        &= \langle c,\, \mathcal{A}(lst\texttt{[$\mathcal{A}(i_1, \sigma, s)$]}\dots\texttt{[$\mathcal{A}(i_d, \sigma, s)$]}, \sigma, s) : e,\, \sigma,\, s,\, pc+1 \rangle 
        && \text{(list flattening)} \\
        &= \langle c,\, \mathcal{A}(lst\texttt{[$i_1$]}\dots\texttt{[$i_d$]}, \sigma, s) : e,\, \sigma,\, s,\, pc+1 \rangle
        && \text{(definition of $\mathcal{A}$)} \\
        &= \langle c,\, \mathcal{A}(a, \sigma, s) : e,\, \sigma,\, s,\, pc+1 \rangle \\
      \end{aligned}
  \]

    \item \textbf{case 6: } $a = \texttt{len($x$)}$ \\
    This case is trivially true as $\texttt{LEN}\;x$ exactly pushes the valuation of \texttt{len($x$)} to the stack.

    \item \textbf{case 7: } $a = \texttt{read()}$ \\
    This case is trivially true as $\texttt{READ}$ exactly pushes the valuation of first unread input to the stack.

    \item \textbf{case 8: } $a = a_1 + a_2$ for some $a_1, a_2 \in \textbf{AExp}$ \\

    Assume for an induction on the correctness of the compilation of $a_1$ and $a_2$.

    \[
      \begin{aligned}
        \langle  \mathcal{CA}(a) : c,\, e,\, \sigma,\, s,\, pc \rangle 
        &= \langle \mathcal{CA}(a_1 + a_2) : c,\, e,\, \sigma,\, s,\, pc \rangle \\
        &= \langle \mathcal{CA}(a_2) : \mathcal{CA}(a_1) : \texttt{ADD} :  c,\, e,\, \sigma,\, s,\, pc \rangle 
        && \text{(definition of $\mathcal{CA}$)} \\
        &\triangleright^{\ast} \langle \mathcal{CA}(a_1) : \texttt{ADD} :  c,\, \mathcal{A}(a_2, \sigma, s) : e,\, \sigma,\, s,\, pc + l_2 \rangle 
        && \text{(induction hypothesis 2)} \\
        &\triangleright^{\ast} \langle \texttt{ADD} :  c,\, \mathcal{A}(a_1, \sigma, s) : \mathcal{A}(a_2, \sigma, s) : e,\, \sigma,\, s,\, pc + l_2 + l_1 \rangle 
        && \text{(induction hypothesis 1)} \\
        &\triangleright \langle c,\, (\mathcal{A}(a_1, \sigma, s) + \mathcal{A}(a_2, \sigma, s)) : e,\, \sigma,\, s,\, pc + l_2 + l_1 + 1 \rangle 
        && \text{(operational semantics)} \\
        &= \langle c,\, \mathcal{A}(a_1 + a_2, \sigma, s) : e,\, \sigma,\, s,\, pc + l_2 + l_1 + 1 \rangle
        && \text{(definition of $\mathcal{A}$)} \\
        &= \langle c,\, \mathcal{A}(a, \sigma, s) : e,\, \sigma,\, s,\, pc + l_2 + l_1 + 1 \rangle \\
      \end{aligned}
  \]

    \item \textbf{case 9: } $a = a_1 \ast a_2$ for some $a_1, a_2 \in \textbf{AExp}$

    Assume for an induction on the correctness of the compilation of $a_1$ and $a_2$.

    \[
      \begin{aligned}
        \langle  \mathcal{CA}(a) : c,\, e,\, \sigma,\, s,\, pc \rangle 
        &= \langle \mathcal{CA}(a_1 \ast a_2) : c,\, e,\, \sigma,\, s,\, pc \rangle \\
        &= \langle \mathcal{CA}(a_2) : \mathcal{CA}(a_1) : \texttt{MUL} :  c,\, e,\, \sigma,\, s,\, pc \rangle 
        && \text{(definition of $\mathcal{CA}$)} \\
        &\triangleright^{\ast} \langle \mathcal{CA}(a_1) : \texttt{MUL} :  c,\, \mathcal{A}(a_2, \sigma, s) : e,\, \sigma,\, s,\, pc + l_2 \rangle 
        && \text{(induction hypothesis 2)} \\
        &\triangleright^{\ast} \langle \texttt{MUL} :  c,\, \mathcal{A}(a_1, \sigma, s) : \mathcal{A}(a_2, \sigma, s) : e,\, \sigma,\, s,\, pc + l_2 + l_1 \rangle 
        && \text{(induction hypothesis 1)} \\
        &\triangleright \langle c,\, (\mathcal{A}(a_1, \sigma, s) \cdot \mathcal{A}(a_2, \sigma, s)) : e,\, \sigma,\, s,\, pc + l_2 + l_1 + 1 \rangle 
        && \text{(operational semantics)} \\
        &= \langle c,\, \mathcal{A}(a_1 \ast a_2, \sigma, s) : e,\, \sigma,\, s,\, pc + l_2 + l_1 + 1 \rangle
        && \text{(definition of $\mathcal{A}$)} \\
        &= \langle c,\, \mathcal{A}(a, \sigma, s) : e,\, \sigma,\, s,\, pc + l_2 + l_1 + 1 \rangle \\
      \end{aligned}
  \]

    \newpage
    \item \textbf{case 10: } $a = a_1 - a_2$ for some $a_1, a_2 \in \textbf{AExp}$

    Assume for an induction on the correctness of the compilation of $a_1$ and $a_2$.

    \[
      \begin{aligned}
        \langle  \mathcal{CA}(a) : c,\, e,\, \sigma,\, s,\, pc \rangle 
        &= \langle \mathcal{CA}(a_1 - a_2) : c,\, e,\, \sigma,\, s,\, pc \rangle \\
        &= \langle \mathcal{CA}(a_2) : \mathcal{CA}(a_1) : \texttt{SUB} :  c,\, e,\, \sigma,\, s,\, pc \rangle 
        && \text{(definition of $\mathcal{CA}$)} \\
        &\triangleright^{\ast} \langle \mathcal{CA}(a_1) : \texttt{SUB} :  c,\, \mathcal{A}(a_2, \sigma, s) : e,\, \sigma,\, s,\, pc + l_2 \rangle 
        && \text{(induction hypothesis 2)} \\
        &\triangleright^{\ast} \langle \texttt{SUB} :  c,\, \mathcal{A}(a_1, \sigma, s) : \mathcal{A}(a_2, \sigma, s) : e,\, \sigma,\, s,\, pc + l_2 + l_1 \rangle 
        && \text{(induction hypothesis 1)} \\
        &\triangleright \langle c,\, (\mathcal{A}(a_1, \sigma, s) - \mathcal{A}(a_2, \sigma, s)) : e,\, \sigma,\, s,\, pc + l_2 + l_1 + 1 \rangle 
        && \text{(operational semantics)} \\
        &= \langle c,\, \mathcal{A}(a_1 - a_2, \sigma, s) : e,\, \sigma,\, s,\, pc + l_2 + l_1 + 1 \rangle
        && \text{(definition of $\mathcal{A}$)} \\
        &= \langle c,\, \mathcal{A}(a, \sigma, s) : e,\, \sigma,\, s,\, pc + l_2 + l_1 + 1 \rangle \\
      \end{aligned}
  \]

    \item \textbf{case 11.1: } $a = a_1 / a_2$ for some $a_1, a_2 \in \textbf{AExp}$ and $(a_1 / a_2)$ has no subterm of type \textbf{Float}

    Assume for an induction on the correctness of the compilation of $a_1$ and $a_2$.

        \begin{itemize}
            \item \textbf{subcase 11.1.1: } $\mathcal{A}(a_2, \sigma, s) \neq 0$

                \[
              \begin{aligned}
                \langle  \mathcal{CA}(a) : c,\, e,\, \sigma,\, s,\, pc \rangle 
                &= \langle \mathcal{CA}(a_1 / a_2) : c,\, e,\, \sigma,\, s,\, pc \rangle \\
                &= \langle \mathcal{CA}(a_2) : \mathcal{CA}(a_1) : \texttt{iDIV} :  c,\, e,\, \sigma,\, s,\, pc \rangle 
                && \text{(definition of $\mathcal{CA}$)} \\
                &\triangleright^{\ast} \langle \mathcal{CA}(a_1) : \texttt{iDIV} :  c,\, \mathcal{A}(a_2, \sigma, s) : e,\, \sigma,\, s,\, pc + l_2 \rangle 
                && \text{(induction hypothesis 2)} \\
                &\triangleright^{\ast} \langle \texttt{iDIV} :  c,\, \mathcal{A}(a_1, \sigma, s) : \mathcal{A}(a_2, \sigma, s) : e,\, \sigma,\, s,\, pc + l_2 + l_1 \rangle 
                && \text{(induction hypothesis 1)} \\
                &\triangleright \langle c,\, \left\lfloor  \mathcal{A} (a_1, \sigma, s) /  \mathcal{A}(a_2, \sigma, s)\right\rfloor : e,\, \sigma,\, s,\, pc + l_2 + l_1 + 1 \rangle 
                && \text{(operational semantics)} \\
                &= \langle c,\, \mathcal{A}(a_1 / a_2, \sigma, s) : e,\, \sigma,\, s,\, pc + l_2 + l_1 + 1 \rangle
                && \text{(definition of $\mathcal{A}$)} \\
                &= \langle c,\, \mathcal{A}(a, \sigma, s) : e,\, \sigma,\, s,\, pc + l_2 + l_1 + 1 \rangle \\
              \end{aligned}
          \]

            \item \textbf{subcase 11.1.2: } $\mathcal{A}(a_2, \sigma, s) = 0$

            \[
              \begin{aligned}
                \langle  \mathcal{CA}(a) : c,\, e,\, \sigma,\, s,\, pc \rangle 
                &= \langle \mathcal{CA}(a_1 / a_2) : c,\, e,\, \sigma,\, s,\, pc \rangle \\
                &= \langle \mathcal{CA}(a_2) : \mathcal{CA}(a_1) : \texttt{iDIV} :  c,\, e,\, \sigma,\, s,\, pc \rangle 
                && \text{(definition of $\mathcal{CA}$)} \\
                &\triangleright^{\ast} \langle \mathcal{CA}(a_1) : \texttt{iDIV} :  c,\, \mathcal{A}(a_2, \sigma, s) : e,\, \sigma,\, s,\, pc + l_2 \rangle 
                && \text{(induction hypothesis 2)} \\
                &= \langle \mathcal{CA}(a_1) : \texttt{iDIV} :  c,\, 0 : e,\, \sigma,\, s,\, pc + l_2 \rangle \\
                &\triangleright^{\ast} \langle \texttt{iDIV} :  c,\, \mathcal{A}(a_1, \sigma, s) : 0 : e,\, \sigma,\, s,\, pc + l_2 + l_1 \rangle 
                && \text{(induction hypothesis 1)} \\
                &\triangleright \langle c,\, \bot : e,\, \hat{\sigma},\, s,\, pc + l_2 + l_1 + 1 \rangle 
                && \text{(operational semantics)} \\
                &= \langle c,\, \mathcal{A}(a_1 / a_2, \sigma, s) : e,\, \hat{\sigma},\, s,\, pc + l_2 + l_1 + 1 \rangle
                && \text{(definition of $\mathcal{A}$)} \\
                &= \langle c,\, \mathcal{A}(a, \sigma, s) : e,\, \hat{\sigma},\, s,\, pc + l_2 + l_1 + 1 \rangle \\
              \end{aligned}
          \]
            
        \end{itemize}

    \item \textbf{case 11.2: } $a = a_1 / a_2$ for some $a_1, a_2 \in \textbf{AExp}$ and $(a_1 / a_2)$ has at least one subterm of type \textbf{Float}

    Assume for an induction on the correctness of the compilation of $a_1$ and $a_2$.

        \begin{itemize}
            \item \textbf{subcase 11.2.1: } $\mathcal{A}(a_2, \sigma, s) \neq 0.0$

                \[
              \begin{aligned}
                \langle  \mathcal{CA}(a) : c,\, e,\, \sigma,\, s,\, pc \rangle 
                &= \langle \mathcal{CA}(a_1 / a_2) : c,\, e,\, \sigma,\, s,\, pc \rangle \\
                &= \langle \mathcal{CA}(a_2) : \mathcal{CA}(a_1) : \texttt{fDIV} :  c,\, e,\, \sigma,\, s,\, pc \rangle 
                && \text{(definition of $\mathcal{CA}$)} \\
                &\triangleright^{\ast} \langle \mathcal{CA}(a_1) : \texttt{fDIV} :  c,\, \mathcal{A}(a_2, \sigma, s) : e,\, \sigma,\, s,\, pc + l_2 \rangle 
                && \text{(induction hypothesis 2)} \\
                &\triangleright^{\ast} \langle \texttt{fDIV} :  c,\, \mathcal{A}(a_1, \sigma, s) : \mathcal{A}(a_2, \sigma, s) : e,\, \sigma,\, s,\, pc + l_2 + l_1 \rangle 
                && \text{(induction hypothesis 1)} \\
                &\triangleright \langle c,\, \mathcal{A} (a_1, \sigma, s) /  \mathcal{A}(a_2, \sigma, s) : e,\, \sigma,\, s,\, pc + l_2 + l_1 + 1 \rangle 
                && \text{(operational semantics)} \\
                &= \langle c,\, \mathcal{A}(a_1 / a_2, \sigma, s) : e,\, \sigma,\, s,\, pc + l_2 + l_1 + 1 \rangle
                && \text{(definition of $\mathcal{A}$)} \\
                &= \langle c,\, \mathcal{A}(a, \sigma, s) : e,\, \sigma,\, s,\, pc + l_2 + l_1 + 1 \rangle \\
              \end{aligned}
          \]

            \item \textbf{subcase 11.2.2: } $\mathcal{A}(a_2, \sigma, s) = 0.0$

            \[
              \begin{aligned}
                \langle  \mathcal{CA}(a) : c,\, e,\, \sigma,\, s,\, pc \rangle 
                &= \langle \mathcal{CA}(a_1 / a_2) : c,\, e,\, \sigma,\, s,\, pc \rangle \\
                &= \langle \mathcal{CA}(a_2) : \mathcal{CA}(a_1) : \texttt{fDIV} :  c,\, e,\, \sigma,\, s,\, pc \rangle 
                && \text{(definition of $\mathcal{CA}$)} \\
                &\triangleright^{\ast} \langle \mathcal{CA}(a_1) : \texttt{fDIV} :  c,\, \mathcal{A}(a_2, \sigma, s) : e,\, \sigma,\, s,\, pc + l_2 \rangle 
                && \text{(induction hypothesis 2)} \\
                &= \langle \mathcal{CA}(a_1) : \texttt{fDIV} :  c,\, 0.0 : e,\, \sigma,\, s,\, pc + l_2 \rangle \\
                &\triangleright^{\ast} \langle \texttt{fDIV} :  c,\, \mathcal{A}(a_1, \sigma, s) : 0.0 : e,\, \sigma,\, s,\, pc + l_2 + l_1 \rangle 
                && \text{(induction hypothesis 1)} \\
                &\triangleright \langle c,\, \bot : e,\, \hat{\sigma},\, s,\, pc + l_2 + l_1 + 1 \rangle 
                && \text{(operational semantics)} \\
                &= \langle c,\, \mathcal{A}(a_1 / a_2, \sigma, s) : e,\, \hat{\sigma},\, s,\, pc + l_2 + l_1 + 1 \rangle
                && \text{(definition of $\mathcal{A}$)} \\
                &= \langle c,\, \mathcal{A}(a, \sigma, s) : e,\, \hat{\sigma},\, s,\, pc + l_2 + l_1 + 1 \rangle \\
              \end{aligned}
          \]
            
        \end{itemize}    
\end{itemize}

where $l_1$ and $l_2$ are the number of commands in $\mathcal{CA}(a_1)$ and $\mathcal{CA}(a_2)$, respectively.\\
Observe that, at each case, the increase in $pc$ is exactly equal to the number of commands in $\mathcal{CA}(a)$, as desired.

\end{proof}

\subsubsection{Boolean}
Given any program stack $e$, non-erroneous program state $\sigma$, scope $s$, program counter $pc$; if $b$ does not have a erroneous subterm:
\[
\forall b \in \textbf{BExp} \,\cdot\, \langle \mathcal{CB}(b) : c,\, e,\, \sigma,\, s,\, pc \rangle \; \triangleright^{\ast} \;
\langle c,\, \mathcal{B}(b, \sigma, s) : e,\, \sigma,\, s,\, pc + l \rangle
\] where $l$ is the number of commands in $\mathcal{CB}(b)$.

\begin{proof}
We will prove this by an induction on the structure of boolean expressions.

\begin{itemize}
    \item \textbf{case 1: } $b = \texttt{true}$

    \[
      \begin{aligned}
        \langle  \mathcal{CB}(b) : c,\, e,\, \sigma,\, s,\, pc \rangle 
        &= \langle \mathcal{CB}(\texttt{true}) : c,\, e,\, \sigma,\, s,\, pc \rangle \\
        &= \langle \texttt{PUSH 1}   : c,\, e,\, \sigma,\, s,\, pc \rangle 
        && \text{(definition of $\mathcal{CB}$)} \\
        &\triangleright \langle c,\, \mathcal{V}(\texttt{1}) : e,\, \sigma,\, s,\, pc+1 \rangle 
        && \text{(operational semantics)} \\
        &= \langle c,\, 1 : e,\, \sigma,\, s,\, pc+1 \rangle  
        \\
        &= \langle c,\, \mathcal{B}(\texttt{true}, \sigma, s) : e,\, \sigma,\, s,\, pc+1 \rangle 
        && \text{(definition of $\mathcal{B}$)} \\
        &= \langle c,\, \mathcal{B}(b, \sigma, s) : e,\, \sigma,\, s,\, pc+1 \rangle \\
      \end{aligned}
  \]

    \item \textbf{case 2: } $b = \texttt{false}$

    \[
      \begin{aligned}
        \langle  \mathcal{CB}(b) : c,\, e,\, \sigma,\, s,\, pc \rangle 
        &= \langle \mathcal{CB}(\texttt{false}) : c,\, e,\, \sigma,\, s,\, pc \rangle \\
        &= \langle \texttt{PUSH 0}   : c,\, e,\, \sigma,\, s,\, pc \rangle 
        && \text{(definition of $\mathcal{CB}$)} \\
        &\triangleright \langle c,\, \mathcal{V}(\texttt{0}) : e,\, \sigma,\, s,\, pc+1 \rangle 
        && \text{(operational semantics)} \\
        &= \langle c,\, 0 : e,\, \sigma,\, s,\, pc+1 \rangle 
        \\
        &= \langle c,\, \mathcal{B}(\texttt{false}, \sigma, s) : e,\, \sigma,\, s,\, pc+1 \rangle 
        && \text{(definition of $\mathcal{B}$)} \\
        &= \langle c,\, \mathcal{B}(b, \sigma, s) : e,\, \sigma,\, s,\, pc+1 \rangle \\
      \end{aligned}
  \]
  
    \newpage
    \item \textbf{case 3: } $b \in \textbf{AExp}$

     \begin{itemize}
         \item \textbf{subcase 3.1: } $\mathcal{A}(a, \sigma, s) \neq 0$

         \[
          \begin{aligned}
            \langle  \mathcal{CB}(b) : c,\, e,\, \sigma,\, s,\, pc \rangle
            &= \langle \mathcal{CB}(a) : c,\, e,\, \sigma,\, s,\, pc \rangle \\
            &= \langle \mathcal{CA}(a) : \texttt{PUSH 0} : \texttt{EQ} : \texttt{NEG} : c,\, e,\, \sigma,\, s,\, pc \rangle
            && \text{(definition of $\mathcal{CB}$)} \\
            &\triangleright^\ast \langle \texttt{PUSH 0} : \texttt{EQ} : \texttt{NEG} : c,\, \mathcal{A}(a, \sigma, s): e,\, \sigma,\, s,\, pc + l_a \rangle
            && \text{(correctness of $\mathcal{CA}$)} \\
            &\triangleright \langle \texttt{EQ} : \texttt{NEG} :  c,\, 0 : \mathcal{A}(a, \sigma, s): e,\, \sigma,\, s,\, pc + l_a + 1 \rangle
            && \text{(operational semantics)} \\
            &\triangleright \langle \texttt{NEG} : c,\, (0 = \mathcal{A}(a, \sigma, s)) : e,\, \sigma,\, s,\, pc + l_a + 2 \rangle
            && \text{(operational semantics)} \\
            &= \langle \texttt{NEG} : c,\, 0 : e,\, \sigma,\, s,\, pc + l_a + 2 \rangle
            && \text{(boolean algebra)} \\
            &\triangleright \langle c,\, 1 : e,\, \sigma,\, s,\, pc + l_a + 3 \rangle
            && \text{(operational semantics)} \\
            &= \langle c,\, \mathcal{B}(a, \sigma, s) : e,\, \sigma,\, s,\, pc + l_a + 3 \rangle
            && \text{(definition of $\mathcal{B}$)} \\
            &= \langle c,\, \mathcal{B}(b, \sigma, s) : e,\, \sigma,\, s,\, pc + l_a + 3 \rangle
            \\
          \end{aligned}
      \]

         \item \textbf{subcase 3.2: } $\mathcal{A}(a, \sigma, s) = 0$

         \[
          \begin{aligned}
            \langle  \mathcal{CB}(b) : c,\, e,\, \sigma,\, s,\, pc \rangle
            &= \langle \mathcal{CB}(a) : c,\, e,\, \sigma,\, s,\, pc \rangle \\
            &= \langle \mathcal{CA}(a) : \texttt{PUSH 0} : \texttt{EQ} : \texttt{NEG} : c,\, e,\, \sigma,\, s,\, pc \rangle
            && \text{(definition of $\mathcal{CB}$)} \\
            &\triangleright^\ast \langle \texttt{PUSH 0} : \texttt{EQ} : \texttt{NEG} : c,\, \mathcal{A}(a, \sigma, s): e,\, \sigma,\, s,\, pc + l_a \rangle
            && \text{(correctness of $\mathcal{CA}$)} \\
            &\triangleright \langle \texttt{EQ} : \texttt{NEG} :  c,\, 0 : \mathcal{A}(a, \sigma, s): e,\, \sigma,\, s,\, pc + l_a + 1 \rangle
            && \text{(operational semantics)} \\
            &\triangleright \langle \texttt{NEG} : c,\, (0 = \mathcal{A}(a, \sigma, s)) : e,\, \sigma,\, s,\, pc + l_a + 2 \rangle
            && \text{(operational semantics)} \\
            &= \langle \texttt{NEG} : c,\, (0 = 0) : e,\, \sigma,\, s,\, pc + l_a + 2 \rangle
            \\
            &= \langle \texttt{NEG} : c,\, 1 : e,\, \sigma,\, s,\, pc + l_a + 2 \rangle
            && \text{(boolean algebra)} \\
            &\triangleright \langle c,\, 0 : e,\, \sigma,\, s,\, pc + l_a + 3 \rangle
            && \text{(operational semantics)} \\
            &= \langle c,\, \mathcal{B}(a, \sigma, s) : e,\, \sigma,\, s,\, pc + l_a + 3 \rangle
            && \text{(definition of $\mathcal{B}$)} \\
            &= \langle c,\, \mathcal{B}(b, \sigma, s) : e,\, \sigma,\, s,\, pc + l_a + 3 \rangle
            \\
          \end{aligned}
      \]
         
     \end{itemize}

    

    \item \textbf{case 4: } $b = (a_1 \;\texttt{=}\; a_2)$ for some $a_1, a_2 \in \textbf{AExp}$

    \begin{itemize}
        \item \textbf{subcase 4.1: } $\mathcal{A}(a_1, \sigma, s) = \mathcal{A}(a_2, \sigma, s)$

        \[
          \begin{aligned}
            \langle  \mathcal{CB}(b) : c,\, e,\, \sigma,\, s,\, pc \rangle
            &= \langle \mathcal{CB}(a_1 \;\texttt{=}\; a_2) : c,\, e,\, \sigma,\, s,\, pc \rangle \\
            &= \langle \mathcal{CA}(a_2) : \mathcal{CA}(a_1) : \texttt{EQ} : c,\, e,\, \sigma,\, s,\, pc \rangle
            && \text{(definition of $\mathcal{CB}$)} \\
            &\triangleright^\ast \langle \mathcal{CA}(a_1) : \texttt{EQ} : c,\, \mathcal{A}(a_2, \sigma, s): e,\, \sigma,\, s,\, pc + l_{a_1} \rangle
            && \text{(correctness of $\mathcal{CA}$)} \\
            &\triangleright^\ast \langle \texttt{EQ} : c,\, \mathcal{A}(a_1, \sigma, s) :\mathcal{A}(a_2, \sigma, s): e,\, \sigma,\, s,\, pc + l_{a_1} + l_{a_2} \rangle
            && \text{(correctness of $\mathcal{CA}$)} \\
            &\triangleright \langle c,\, (\mathcal{A}(a_1, \sigma, s) = \mathcal{A}(a_2, \sigma, s)) : e,\, \sigma,\, s,\, pc + l_{a_1} + l_{a_2} + 1 \rangle
            && \text{(operational semantics)} \\
            &= \langle c,\, 1 : e,\, \sigma,\, s,\, pc + l_{a_1} + l_{a_2} + 1 \rangle \\
            &= \langle c,\, \mathcal{B}(a_1 \;\texttt{=}\; a_2, \sigma, s) : e,\, \sigma,\, s,\, pc + l_{a_1} + l_{a_2} + 1 \rangle
            && \text{(definition of $\mathcal{B}$)} \\
            &= \langle c,\, \mathcal{B}(b, \sigma, s) : e,\, \sigma,\, s,\, pc + l_{a_1} + l_{a_2} + 1\rangle
            \\
          \end{aligned}
      \]

        \item \textbf{subcase 4.1: } $\mathcal{A}(a_1, \sigma, s) \neq \mathcal{A}(a_2, \sigma, s)$

        \[
          \begin{aligned}
            \langle  \mathcal{CB}(b) : c,\, e,\, \sigma,\, s,\, pc \rangle
            &= \langle \mathcal{CB}(a_1 \;\texttt{=}\; a_2) : c,\, e,\, \sigma,\, s,\, pc \rangle \\
            &= \langle \mathcal{CA}(a_2) : \mathcal{CA}(a_1) : \texttt{EQ} : c,\, e,\, \sigma,\, s,\, pc \rangle
            && \text{(definition of $\mathcal{CB}$)} \\
            &\triangleright^\ast \langle \mathcal{CA}(a_1) : \texttt{EQ} : c,\, \mathcal{A}(a_2, \sigma, s): e,\, \sigma,\, s,\, pc + l_{a_1} \rangle
            && \text{(correctness of $\mathcal{CA}$)} \\
            &\triangleright^\ast \langle \texttt{EQ} : c,\, \mathcal{A}(a_1, \sigma, s) :\mathcal{A}(a_2, \sigma, s): e,\, \sigma,\, s,\, pc + l_{a_1} + l_{a_2} \rangle
            && \text{(correctness of $\mathcal{CA}$)} \\
            &\triangleright \langle c,\, (\mathcal{A}(a_1, \sigma, s) = \mathcal{A}(a_2, \sigma, s)) : e,\, \sigma,\, s,\, pc + l_{a_1} + l_{a_2} + 1 \rangle
            && \text{(operational semantics)} \\
            &= \langle c,\, 0 : e,\, \sigma,\, s,\, pc + l_{a_1} + l_{a_2} + 1 \rangle \\
            &= \langle c,\, \mathcal{B}(a_1 \;\texttt{=}\; a_2, \sigma, s) : e,\, \sigma,\, s,\, pc + l_{a_1} + l_{a_2} + 1 \rangle
            && \text{(definition of $\mathcal{B}$)} \\
            &= \langle c,\, \mathcal{B}(b, \sigma, s) : e,\, \sigma,\, s,\, pc + l_{a_1} + l_{a_2} + 1 \rangle
            \\
          \end{aligned}
      \]
    \end{itemize}
    

    \item \textbf{case 5: } $b = (\texttt{not}\; b')$ for some $b' \in \textbf{BExp}$

    Assume for an induction on the correctness of the compilation of $b'$.

    \[
      \begin{aligned}
        \langle  \mathcal{CB}(b) : c,\, e,\, \sigma,\, s,\, pc \rangle 
        &= \langle \mathcal{CB}(\texttt{not}\; b') : c,\, e,\, \sigma,\, s,\, pc \rangle \\
        &= \langle \mathcal{CB}(b') : \texttt{NEG} : c,\, e,\, \sigma,\, s,\, pc \rangle 
        && \text{(definition of $\mathcal{CB}$)} \\
        &\triangleright^\ast \langle \texttt{NEG} : c,\, \mathcal{B}(b', \sigma, s) : e,\, \sigma,\, s,\, pc + l_{b'} \rangle 
        && \text{(induction hypothesis)} \\
        &\triangleright \langle c,\, \neg\mathcal{B}(b', \sigma, s) : e,\, \sigma,\, s,\, pc + l_{b'} + 1 \rangle 
        && \text{(operational semantics)} \\
        &= \langle c,\, \mathcal{B}(\texttt{not}\; b', \sigma, s) : e,\, \sigma,\, s,\, pc + l_{b'} + 1 \rangle 
        && \text{(definition of $\mathcal{B}$)} \\
        &= \langle c,\, \mathcal{B}(b, \sigma, s) : e,\, \sigma,\, s,\, pc + l_{b'} + 1 \rangle \\
      \end{aligned}
  \]

    \item \textbf{case 6: } $b = (b_1 \;\texttt{and}\; b_2)$ for some $b_1, b_2 \in \textbf{BExp}$

    Assume for an induction on the correctness of the compilations of $b_1$ and $b_2$.

    \[
      \begin{aligned}
        \langle  \mathcal{CB}(b) : c,\, e,\, \sigma,\, s,\, pc \rangle 
        &= \langle \mathcal{CB}(b_1 \;\texttt{and}\; b_2) : c,\, e,\, \sigma,\, s,\, pc \rangle \\
        &= \langle \mathcal{CB}(b_2) : \mathcal{CB}(b_1) : \texttt{AND} : c,\, e,\, \sigma,\, s,\, pc \rangle 
        && \text{(definition of $\mathcal{CB}$)} \\
        &\triangleright^\ast \langle \mathcal{CB}(b_1) : \texttt{AND} : c,\, \mathcal{B}(b_2, \sigma, s) : e,\, \sigma,\, s,\, pc + l_{b_2} \rangle 
        && \text{(induction hypothesis 2)} \\
        &\triangleright^\ast \langle \texttt{AND} : c,\, \mathcal{B}(b_1, \sigma, s) : \mathcal{B}(b_2, \sigma, s) : e,\, \sigma,\, s,\, pc + l_{b_2} + l_{b_1} \rangle 
        && \text{(induction hypothesis 2)} \\
        &\triangleright \langle c,\, (\mathcal{B}(b_1, \sigma, s) \,\land\, \mathcal{B}(b_2, \sigma, s) : e,\, \sigma,\, s,\, pc + l_{b_2} + l_{b_1} + 1 \rangle
        && \text{(operational semantics)} \\
        &= \langle c,\, \mathcal{B}(b_1 \;\texttt{and}\; b_2, \sigma, s) : e,\, \sigma,\, s,\, pc + l_{b_2} + l_{b_1} + 1 \rangle 
        && \text{(definition of $\mathcal{B}$)} \\
        &= \langle c,\, \mathcal{B}(b, \sigma, s) : e,\, \sigma,\, s,\, pc + l_{b_2} + l_{b_1} + 1 \rangle \\
      \end{aligned}
  \]
\end{itemize}

where $l_{a'}$ denotes the number of commands in $\mathcal{CA}(a')$ for all $a' \in \textbf{AExp}$ and $l_{b'}$ denotes the number of commands in $\mathcal{CB}(b')$ for all $b' \in \textbf{BExp}$.\\
Observe that in each case the increase in $pc$ is exactly equal to the number of commands in $\mathcal{CB}(b)$, as desired.

\end{proof}

\subsection{Correctness of the statements}

Now, as the correctness of arithmetic and boolean expressions is proven, we may restate and prove the correctness of our compiler:

\[
\langle P,\, \underline{\sigma},\, s \rangle \; \rightarrow^{\ast} \; \langle \underline{\sigma}_f,\, s_f \rangle \quad \implies \quad 
\langle \mathcal{C}(P),\, e,\, \underline{\sigma},\, s,\, pc,\, ras \rangle \; \triangleright^{\ast} \; \langle \underline{\sigma}_f,\, s_f,\, pc + l_P,\, ras \rangle 
\] where $P$ is any \texttt{SelVeri} program and $l_P$ is the number of commands in $\mathcal{C}(P)$. \\

Observe that the implication above mentions the correctness of only terminating programs. However, we also claim that the compilation of non-terminating \texttt{SelVeri} programs are also correct, and we will present our arguments in \textbf{section 6}.

\begin{proof}
If $\underline{\sigma}$ is an erroneous program state, then by section \textbf{2.2.5} of this document and the operational semantic rule of high-level \texttt{SelVeri} language $[ \text{err}_{\text{os}} ]$, the above implication is obviously true \footnote{Furthermore, $\sigma = \sigma_f$ and $s = s_f$, in this case.}.  Therefore, from now on, we assume that $\underline{\sigma}$ represents a non-erroneous program state and will use $\sigma$ to denote it. \\

We will prove the correctness of $\mathcal{C}$ via a strong induction on the number of execution steps, by utilizing the fact that sub-statements naturally require a lesser number of execution steps. The correctness of basic type declarations, I/O and \texttt{pass} statements establish our base cases, since their compilation include only one command.

\newpage
\begin{itemize}
    \item \textbf{(base) case 1:} $P = (x \;\texttt{:}\; t)$

        \begin{itemize}
            \item \textbf{subcase 1.1:} $t \in \textbf{BasicType} \land s(x) = \textbf{Void}$

            \[
              \begin{aligned}
              \langle P,\, \sigma,\, s \rangle
              &=\langle x \;\texttt{:}\; t,\, \sigma,\, s \rangle \\
              &\rightarrow \langle \sigma[x \mapsto \textbf{0}] ,\, s[x \mapsto t] \rangle
              && [ \text{decl}_{\text{os}} ] \\   
              \end{aligned}
          \]
          
          and

          \[
              \begin{aligned}
              \langle \mathcal{C}(P),\, e,\, \sigma,\, s,\, pc \rangle
              &= \langle \mathcal{C}(x \;\texttt{:}\; t),\, e,\, \sigma,\, s,\, pc \rangle \\
              &= \langle \texttt{DECL t},\, e,\, \sigma,\, s,\, pc \rangle
              && \text{(definition of $\mathcal{C}$)} \\   
              &= \langle \epsilon,\, e,\, \sigma[x \mapsto \textbf{0}] ,\, s[x \mapsto t],\, pc+1 \rangle
              && \text{(operational semantics)} \\ 
              \end{aligned}
          \]

            \item \textbf{subcase 1.2:} $t \notin \textbf{BasicType}  \lor s(x) \neq \textbf{Void}$

            \[
              \begin{aligned}
              \langle P,\, \sigma,\, s \rangle
              &=\langle x \;\texttt{:}\; t,\, \sigma,\, s \rangle \\
              &\rightarrow \hat{\sigma}
              && [ \text{decl}^\bot_{\text{os}} ] \\   
              \end{aligned}
          \]

          and

          \[
              \begin{aligned}
              \langle \mathcal{C}(P),\, e,\, \sigma,\, s,\, pc \rangle
              &= \langle \mathcal{C}(x \;\texttt{:}\; t),\, e,\, \sigma,\, s,\, pc \rangle \\
              &= \langle \texttt{DECL t},\, e,\, \sigma,\, s,\, pc \rangle
              && \text{(definition of $\mathcal{C}$)} \\   
              &= \langle \epsilon,\, e,\, \hat{\sigma} ,\, s,\, pc+1 \rangle
              && \text{(operational semantics)} \\ 
              \end{aligned}
          \]
            
        \end{itemize}

    \item \textbf{(base) case 2:} $P =\texttt{pass}$\\

    \[
      \begin{aligned}
      \langle P,\, \sigma,\, s \rangle
      &=\langle \texttt{pass},\, \sigma,\, s \rangle \\
      &\rightarrow \langle \sigma ,\, s \rangle
      && [ \text{pass}_{\text{os}} ] \\   
      \end{aligned}
  \]
  
  and

  \[
      \begin{aligned}
      \langle \mathcal{C}(P),\, e,\, \sigma,\, s,\, pc \rangle
      &= \langle \mathcal{C}(\texttt{pass}),\, e,\, \sigma,\, s,\, pc \rangle \\
      &= \langle \texttt{NOOP},\, e,\, \sigma,\, s,\, pc \rangle
      && \text{(definition of $\mathcal{C}$)} \\   
      &= \langle \epsilon,\, e,\, \sigma ,\, s,\, pc+1 \rangle
      && \text{(operational semantics)} \\ 
      \end{aligned}
  \]

    \item \textbf{(base) case 3:} $P$ is an IO-statement, i.e. $P \in \{ \texttt{write}\;x ,\, \texttt{writeline}\; x \}$. Then, the compilation of IO-statement is trivially true as they compile into one IR-command whose semantics apply the exact same change to the program state and scope.

    \item \textbf{case 4:} $P = (lst \;\texttt{:}\; \textbf{List[$t, d, \vec{a}$]})$\\

    Let $\sigma' = \sigma[\texttt{\_$lst$\_len\_i} \mapsto \mathcal{A}(a_i, \sigma, s)]$ and $s' = s[\texttt{\_$lst$\_len\_i} \mapsto \textbf{Int}]$ for $ i \in \{1,\dots,d\}$ 

        \[
          \begin{aligned}
          \langle P,\, \sigma,\, s \rangle
          &=\langle lst \;\texttt{:}\; \textbf{List[$t, d, \vec{a}$]},\, \sigma,\, s \rangle \\
          &\rightarrow \langle \sigma'[lst \mapsto \textbf{0}], s'[lst \mapsto \textbf{List[$t, d, \vec{\mathcal{A}(a_i, \sigma, s)}$]}] \rangle
          && [ \text{decl}^\text{list}_{\text{os}} ]\\   
          \end{aligned}
      \]
    
      and

          \[
              \begin{aligned}
              \langle \mathcal{C}(P),\, e,\, \sigma,\, s,\, pc \rangle
              &= \langle \mathcal{C}(lst \;\texttt{:}\; \textbf{List[$t, a$]}),\, e,\, \sigma,\, s,\, pc \rangle \\
              &= \langle \texttt{DECL INT}\;\texttt{\_$lst$\_len\_1} : \mathcal{CA}(a_1) : \texttt{STORE}\;\texttt{\_$lst$\_len\_1} : \dots : \texttt{DECL INT}\;\texttt{\_$lst$\_len\_$d$} : \mathcal{CA}(a_d) : \texttt{STORE}\;\texttt{\_$lst$\_len\_$d$} : \mathcal{CA}(\texttt{size\_IR(} \textbf{List[$t,d,\vec{a}$]} \texttt{)} ) : \texttt{LDECL LIST[$t$]}\;lst,\, e,\, \sigma,\, s,\, pc \rangle
              && \text{(definition of $\mathcal{C}$)} \\
              &\triangleright^\ast \langle \mathcal{CA}(\texttt{size\_IR(} \textbf{List[$t, a$]} \texttt{)} ) : \texttt{LDECL LIST[$t$]}\;lst,\, e,\, \sigma',\, s',\, pc + l \rangle
              && \text{(operational semantics))} \\
              &\triangleright \langle \texttt{LDECL LIST[$t$]}\;lst,\, \mathcal{A}(\texttt{size\_IR(} \textbf{List[$t, d, \vec{a}$]} \texttt{)}, \sigma, s) : e,\, \sigma[\texttt{\_$lst$\_len\_1} \mapsto \mathcal{A}(a, \sigma, s)],\, s[\texttt{\_$lst$\_len\_1} \mapsto \texttt{INT}],\, pc + l + l_a \rangle
              && \text{(correctness of $\mathcal{CA}$)} \\ 
              &\triangleright \langle \epsilon,\, e,\, \sigma'[lst \mapsto \textbf{0}],\, s'[lst \mapsto \texttt{LIST[$t, \mathcal{A}(\texttt{size\_IR(} \textbf{List[$t, d, \vec{a}$]} \texttt{)}, \sigma, s)$]}],\, pc + l + l_a + 1\rangle
              && \text{(operational semantics)} \\ 
              &= \langle \epsilon,\, e,\, \sigma'[lst \mapsto \textbf{0}],\, s'[lst \mapsto \textbf{List[$t, d, \vec{\mathcal{A}(a_i, \sigma, s)}$]}],\,  pc + l + l_a + 2 \rangle
              && \text{(list flattening)} \\ 
              \end{aligned}
          \]

        Note that, for all use cases of $\mathcal{A}$ in the proof above, the valuation does not change whether $\sigma, s$ or $\sigma', s'$ is used, as only changes between two tuples are related to newly declared variables.
  

    \item \textbf{case 5:} $P = (x \;\texttt{:=}\; a)$\\

    Assume that $x$ is of a basic type and $\mathcal{A}(a, \sigma, s) \neq \bot \,\land\, \mathcal{A}(a, \sigma, s) \in ran\mathcal{L}_{s(x)}$, then:

    \[
      \begin{aligned}
      \langle P,\, \sigma,\, s \rangle
      &=\langle x \;\texttt{:=}\; a,\, \sigma,\, s \rangle \\
      &\sigma[x \mapsto \mathcal{A}(a, \sigma, s)]
      && [ \text{ass}_{\text{os}} ] \\   
      \end{aligned}
  \]
  
  then

  \[
      \begin{aligned}
      \langle \mathcal{C}(P),\, e,\, \sigma,\, s,\, pc \rangle
      &= \langle \mathcal{C}(x \;\texttt{:=}\; a),\, e,\, \sigma,\, s,\, pc \rangle \\
      &= \langle \mathcal{CA}(a) : \texttt{STORE}\;x ,\, e,\, \sigma,\, s,\, pc \rangle
      && \text{(definition of $\mathcal{C}$)} \\   
      &\triangleright^\ast \langle \texttt{STORE}\;x ,\, \mathcal{A}(a, \sigma, s) : e,\, \sigma,\, s,\, pc + l_a \rangle
      && \text{(correctness of $\mathcal{CA}$)} \\ 
      &= \langle \epsilon ,\,  e,\, \sigma[x \mapsto \mathcal{A}(a, \sigma, s)],\, s,\, pc + l_a + 1 \rangle
      && \text{(operational semantics)} \\ 
      \end{aligned}
  \]
  \item \textbf{case 6:} $P = (lst\texttt{[$i_1$]}\dots\texttt{[$i_d$]} \;\texttt{:=}\; a)$\\

    Assume that $lst$ is an $d$-dimensional list and $\mathcal{A}(a, \sigma, s) \neq \bot \,\land\, \mathcal{A}(a, \sigma, s) \in ran\mathcal{L}_{s(x)}$, then:

    \[
      \begin{aligned}
      \langle P,\, \sigma,\, s \rangle
      &=\langle lst\texttt{[$i_1$]}\dots\texttt{[$i_d$]} \;\texttt{:=}\; a,\, \sigma,\, s \rangle \\
      &\sigma[lst\texttt{[$\mathcal{A}(i_1, \sigma, s)$]}\dots\texttt{[$\mathcal{A}(i_d, \sigma, s)$]} \mapsto \mathcal{A}(a, \sigma, s)]
      && [ \text{ass}^{\text{list}}_{\text{os}} ] \\   
      \end{aligned}
  \]
  
  then

  \[
      \begin{aligned}
      \langle \mathcal{C}(P),\, e,\, \sigma,\, s,\, pc \rangle
      &= \langle \mathcal{C}(lst\texttt{[$i_1$]}\dots\texttt{[$i_d$]} \;\texttt{:=}\; a),\, e,\, \sigma,\, s,\, pc \rangle \\
      &= \langle \mathcal{CA}(a) : \mathcal{CA}(\mathit{LF}(lst; i_1, \dots, i_d)) :  \texttt{LSTORE}\;lst ,\, e,\, \sigma,\, s,\, pc \rangle
      && \text{(definition of $\mathcal{C}$)} \\   
      &\triangleright^\ast \langle \mathcal{CA}(\mathit{LF}(lst; i_1, \dots, i_d)) :  \texttt{LSTORE}\;lst ,\, \mathcal{A}(a, \sigma, s) : e,\, \sigma,\, s,\, pc + l_a \rangle
      && \text{(correctness of $\mathcal{CA}$)} \\ 
      &\triangleright^\ast \langle \texttt{LSTORE}\;lst ,\, \mathcal{A}(\mathit{LF}(lst\;; i_1, \dots, i_d), \sigma, s) : \mathcal{A}(a, \sigma, s) : e,\, \sigma,\, s,\, pc + l_a + l_I \rangle
      && \text{(correctness of $\mathcal{CA}$)} \\ 
      &\triangleright \langle \epsilon ,\, e,\, \sigma[lst\texttt{[$\mathcal{A}(\mathit{LF}(lst\;; i_1, \dots, i_d), \sigma, s)$]} \mapsto \mathcal{A}(a, \sigma, s)],\, s,\, pc + l_a + l_I + 1\rangle
      && \text{(operational semantics)} \\
      &= \langle \epsilon ,\, e,\, \sigma[lst\texttt{[$\mathcal{A}(i_1, \sigma, s)$]}\dots\texttt{[$\mathcal{A}(i_d, \sigma, s)$]} \mapsto \mathcal{A}(a, \sigma, s)],\, s,\, pc + l_a + l_I + 1 \rangle
      && \text{(list flattening)} \\ 
      \end{aligned}
  \]

    \item \textbf{case 7:} $P = (\texttt{if } b \texttt{ then } S \texttt{ fi})$\\

        \begin{itemize}
            \item \textbf{subcase 7.1:} $\mathcal{B}(b, \sigma, s) = 1$

            Assume

                \[
                  \begin{aligned}
                  \langle P,\, \sigma,\, s \rangle
                  &=\langle \texttt{if } b \texttt{ then } S \texttt{ fi},\, \sigma,\, s \rangle \\
                  &=\langle \texttt{if } b \texttt{ then } S \texttt{ else pass fi},\, \sigma,\, s \rangle 
                  && \text{(semantical equivalance)}\\
                  &\rightarrow \langle S\texttt{;endscope},\, \tilde{\sigma_c} ,\, \tilde{s_c} \rangle
                  && [\text{if}^1_{\text{os}}]   \\
                  &\rightarrow^\ast \langle \texttt{endscope},\, \sigma'_c ,\, s'_c \rangle 
                  && ([\text{comp}^1_{\text{os}}] \,,\, [\text{comp}^2_{\text{os}}])   \\
                  &\rightarrow \langle \underline{\sigma}' ,\, s' \rangle
                  && [\text{scope}_{\text{os}}]   \\
                  \end{aligned}
              \]

            and assume for an induction hypothesis, that is, the compilation of $S$ is correct, then

            \[
              \begin{aligned}
              \langle \mathcal{C}(P),\, e,\, \sigma,\, s,\, pc \rangle
              &= \langle \mathcal{C}(\texttt{if } b \texttt{ then } S \texttt{ fi}),\, e,\, \sigma,\, s,\, pc \rangle \\
              &= \langle \mathcal{CB}(b) : \texttt{JZ}\; pc_{\texttt{NOOP}} : \texttt{CSCOPE} : \mathcal{C}(S) : \texttt{PSCOPE} :\texttt{NOOP} ,\, e,\, \sigma,\, s,\, pc \rangle
              && \text{(definition of $\mathcal{C}$)} \\   
              &\triangleright^\ast \langle \texttt{JZ}\; pc_{\texttt{NOOP}} : \texttt{CSCOPE} : \mathcal{C}(S) : \texttt{PSCOPE} :\texttt{NOOP} ,\, \mathcal{B}(b, \sigma, s) : e,\, \sigma,\, s,\, pc + l_b \rangle
              && \text{(correctness of $\mathcal{CB}$)} \\ 
              &= \langle \texttt{JZ}\; pc_{\texttt{NOOP}} : \texttt{CSCOPE} : \mathcal{C}(S) : \texttt{PSCOPE} :\texttt{NOOP} ,\, 1 : e,\, \sigma,\, s,\, pc + l_b \rangle \\ 
              &\triangleright \langle \texttt{CSCOPE} : \mathcal{C}(S) : \texttt{PSCOPE} :\texttt{NOOP} ,\, e,\, \sigma,\, s,\, pc + l_b + 1 \rangle
              && \text{(operational semantics)} \\
              &\triangleright \langle \mathcal{C}(S) : \texttt{PSCOPE} : \texttt{NOOP} ,\, e,\, \tilde{\sigma_c},\, \tilde{s_c},\, pc + l_b + 2 \rangle
              && \text{(operational semantics)} \\
              &\triangleright^\ast \langle \texttt{PSCOPE} : \texttt{NOOP} ,\, e',\, \sigma_c',\, s_c',\, pc + l_b + l_S + 2 \rangle
              && \text{(induction hypothesis)} \\
              &\triangleright \langle \texttt{NOOP} ,\, e',\, \underline{\sigma}',\, s',\, pc + l_b + l_S + 3 \rangle
              && \text{(induction hypothesis)} \\
              &\triangleright \langle \epsilon ,\, e',\, \underline{\sigma}',\, s',\, pc + l_b + l_S + 4 \rangle
              && \text{(operational semantics)} \\ 
              \end{aligned}
          \]


        \item \textbf{subcase 7.2:} $\mathcal{B}(b, \sigma, s) = 0$

            Assume

                \[
                  \begin{aligned}
                  \langle P,\, \sigma,\, s \rangle
                  &=\langle \texttt{if } b \texttt{ then } S \texttt{ fi},\, \sigma,\, s \rangle \\
                  &=\langle \texttt{if } b \texttt{ then } S \texttt{ else pass fi},\, \sigma,\, s \rangle 
                  && \text{(semantical equivalance)}\\
                  &\rightarrow \langle \texttt{pass;endscope},\, \tilde{\sigma_c} ,\, \tilde{s_c} \rangle
                  && [\text{if}^0_{\text{os}}]   \\
                  &\rightarrow^\ast \langle \texttt{endscope},\, \tilde{\sigma_c} ,\, \tilde{s_c} \rangle 
                  && ([\text{pass}_{\text{os}}] \,,\, [\text{comp}^1_{\text{os}}] \,,\, [\text{comp}^2_{\text{os}}])   \\
                  &\rightarrow \langle \sigma ,\, s \rangle
                  && [\text{scope}_{\text{os}}]   \\
                  \end{aligned}
              \]

            then,

            \[
              \begin{aligned}
              \langle \mathcal{C}(P),\, e,\, \sigma,\, s,\, pc \rangle
              &= \langle \mathcal{C}(\texttt{if } b \texttt{ then } S \texttt{ fi}),\, e,\, \sigma,\, s,\, pc \rangle \\
              &= \langle \mathcal{CB}(b) : \texttt{JZ}\; pc_{\texttt{NOOP}} : \texttt{CSCOPE} : \mathcal{C}(S) : \texttt{PSCOPE} :\texttt{NOOP} ,\, e,\, \sigma,\, s,\, pc \rangle
              && \text{(definition of $\mathcal{C}$)} \\   
              &\triangleright^\ast \langle \texttt{JZ}\; pc_{\texttt{NOOP}} : \texttt{CSCOPE} : \mathcal{C}(S) : \texttt{PSCOPE} :\texttt{NOOP} ,\, \mathcal{B}(b, \sigma, s) : e,\, \sigma,\, s,\, pc + l_b \rangle
              && \text{(correctness of $\mathcal{CB}$)} \\ 
              &= \langle \texttt{JZ}\; pc_{\texttt{NOOP}} : \texttt{CSCOPE} : \mathcal{C}(S) : \texttt{PSCOPE} :\texttt{NOOP} ,\, 0 : e,\, \sigma,\, s,\, pc + l_b \rangle \\  
              &\triangleright \langle \texttt{NOOP} ,\, e,\, \sigma,\, s,\, pc_{\texttt{NOOP}} \rangle
              && \text{(operational semantics)} \\
              &\triangleright \langle \epsilon ,\, e,\, \sigma,\, s,\ pc_{\texttt{NOOP}} + 1 \rangle
              && \text{(operational semantics)} \\ 
              \end{aligned}
          \]
                
        \end{itemize}

    \item \textbf{case 8:} $P = (\texttt{if } b \texttt{ then } S_1 \texttt{ else } S_2 \texttt{ fi})$\\

    \begin{itemize}
            \item \textbf{subcase 8.1:} $\mathcal{B}(b, \sigma, s) = 1$

            Assume

                \[
                  \begin{aligned}
                  \langle P,\, \sigma,\, s \rangle
                  &=\langle \texttt{if } b \texttt{ then } S_1 \texttt{ else } S_2 \texttt{ fi},\, \sigma,\, s \rangle \\
                  &\rightarrow \langle S_1\texttt{;endscope},\, \tilde{\sigma_c} ,\, \tilde{s_c} \rangle
                  && [\text{if}^1_{\text{os}}]   \\
                  &\rightarrow^\ast \langle \texttt{endscope},\, \sigma'_c ,\, s'_c \rangle 
                  && ([\text{comp}^1_{\text{os}}] \,,\, [\text{comp}^2_{\text{os}}])   \\
                  &\rightarrow \langle \underline{\sigma}' ,\, s' \rangle
                  && [\text{scope}_{\text{os}}]   \\
                  \end{aligned}
              \]

            and assume for an induction hypothesis, that is, the compilation of $S_1$ is correct, then

            \[
              \begin{aligned}
              \langle \mathcal{C}(P),\, e,\, \sigma,\, s,\, pc \rangle
              &= \langle \mathcal{C}(\texttt{if } b \texttt{ then } S_1 \texttt{ else } S_2 \texttt{ fi}),\, e,\, \sigma,\, s,\, pc \rangle \\
              &= \langle \mathcal{CB}(b) : \texttt{JZ}\; pc_{\texttt{CSCOPE}_2} : \texttt{CSCOPE} : \mathcal{C}(S_1) : \texttt{PSCOPE} : \texttt{GOTO}\; pc_{\texttt{NOOP}} : \texttt{CSCOPE} : \mathcal{C}(S_2) : \texttt{PSCOPE} : \texttt{NOOP} ,\, e,\, \sigma,\, s,\, pc \rangle
              && \text{(definition of $\mathcal{C}$)} \\
              &\triangleright^\ast \langle \texttt{JZ}\; pc_{\texttt{CSCOPE}_2} : \texttt{CSCOPE} : \mathcal{C}(S_1) : \texttt{PSCOPE} : \texttt{GOTO}\; pc_{\texttt{NOOP}} : \texttt{CSCOPE} : \mathcal{C}(S_2) : \texttt{PSCOPE} : \texttt{NOOP} ,\, \mathcal{B}(b, \sigma, s) : e,\, \sigma,\, s,\, pc + l_b \rangle
              && \text{(correctness of $\mathcal{CB}$)} \\
              &= \langle \texttt{JZ}\; pc_{\texttt{CSCOPE}_2} : \texttt{CSCOPE} : \mathcal{C}(S_1) : \texttt{PSCOPE} : \texttt{GOTO}\; pc_{\texttt{NOOP}} : \texttt{CSCOPE} : \mathcal{C}(S_2) : \texttt{PSCOPE} : \texttt{NOOP},\, 1 : e,\, \sigma,\, s,\, pc + l_b \rangle \\
              &\triangleright \langle \texttt{CSCOPE} : \mathcal{C}(S_1) : \texttt{PSCOPE} : \texttt{GOTO}\; pc_{\texttt{NOOP}} : \texttt{CSCOPE} : \mathcal{C}(S_2) : \texttt{PSCOPE} : \texttt{NOOP},\, e,\, \sigma,\, s,\, pc + l_b + 1 \rangle
              && \text{(operational semantics)} \\
              &\triangleright \langle \mathcal{C}(S_1) : \texttt{PSCOPE} : \texttt{GOTO}\; pc_{\texttt{NOOP}} : \texttt{CSCOPE} : \mathcal{C}(S_2) : \texttt{PSCOPE} : \texttt{NOOP},\, e,\, \tilde{\sigma_c},\, \tilde{s_c},\, pc + l_b + 2 \rangle
              && \text{(operational semantics)} \\
              &\triangleright^\ast \langle \texttt{PSCOPE} : \texttt{GOTO}\; pc_{\texttt{NOOP}} : \texttt{CSCOPE} : \mathcal{C}(S_2) : \texttt{PSCOPE} : \texttt{NOOP},\, e',\, \sigma'_c ,\, s'_c,\, pc + l_b + l_{S_1} + 2 \rangle
              && \text{(induction hypothesis)} \\
              &\triangleright \langle \texttt{GOTO}\; pc_{\texttt{NOOP}} : \texttt{CSCOPE} : \mathcal{C}(S_2) : \texttt{PSCOPE} : \texttt{NOOP},\, e',\, \underline{\sigma}' ,\, s',\, pc + l_b + l_{S_1} + 3 \rangle
              && \text{(induction hypothesis)} \\
              &\triangleright \langle \texttt{NOOP},\, e',\, \underline{\sigma}' ,\, s',\, pc_{\texttt{NOOP}} \rangle
              && \text{(operational semantics)} \\
              &\triangleright \langle \epsilon,\, e',\, \underline{\sigma}' ,\, s',\, pc_{\texttt{NOOP}} + 1 \rangle
              && \text{(operational semantics)} \\
              \end{aligned}
          \]

        \item \textbf{subcase 8.2:} $\mathcal{B}(b, \sigma, s) = 0$

            Assume

                \[
                  \begin{aligned}
                  \langle P,\, \sigma,\, s \rangle
                  &=\langle \texttt{if } b \texttt{ then } S_1 \texttt{ else } S_2 \texttt{ fi},\, \sigma,\, s \rangle \\
                  &\rightarrow \langle S_2\texttt{;endscope},\, \tilde{\sigma_c} ,\, \tilde{s_c} \rangle
                  && [\text{if}^1_{\text{os}}]   \\
                  &\rightarrow^\ast \langle \texttt{endscope},\, \sigma'_c ,\, s'_c \rangle 
                  && ([\text{comp}^1_{\text{os}}] \,,\, [\text{comp}^2_{\text{os}}])   \\
                  &\rightarrow \langle \underline{\sigma}' ,\, s' \rangle
                  && [\text{scope}_{\text{os}}]   \\
                  \end{aligned}
              \]

            and assume for an induction hypothesis, that is, the compilation of $S_2$ is correct, then

            \[
              \begin{aligned}
              \langle \mathcal{C}(P),\, e,\, \sigma,\, s,\, pc \rangle
              &= \langle \mathcal{C}(\texttt{if } b \texttt{ then } S_1 \texttt{ else } S_2 \texttt{ fi}),\, e,\, \sigma,\, s,\, pc \rangle \\
              &= \langle \mathcal{CB}(b) : \texttt{JZ}\; pc_{\texttt{CSCOPE}_2} : \texttt{CSCOPE} : \mathcal{C}(S_1) : \texttt{PSCOPE} : \texttt{GOTO}\; pc_{\texttt{NOOP}} : \texttt{CSCOPE} : \mathcal{C}(S_2) : \texttt{PSCOPE} : \texttt{NOOP} ,\, e,\, \sigma,\, s,\, pc \rangle
              && \text{(definition of $\mathcal{C}$)} \\
              &\triangleright^\ast \langle \texttt{JZ}\; pc_{\texttt{CSCOPE}_2} : \texttt{CSCOPE} : \mathcal{C}(S_1) : \texttt{PSCOPE} : \texttt{GOTO}\; pc_{\texttt{NOOP}} : \texttt{CSCOPE} : \mathcal{C}(S_2) : \texttt{PSCOPE} : \texttt{NOOP} ,\, \mathcal{B}(b, \sigma, s) : e,\, \sigma,\, s,\, pc + l_b \rangle
              && \text{(correctness of $\mathcal{CB}$)} \\
              &= \langle \texttt{JZ}\; pc_{\texttt{CSCOPE}_2} : \texttt{CSCOPE} : \mathcal{C}(S_1) : \texttt{PSCOPE} : \texttt{GOTO}\; pc_{\texttt{NOOP}} : \texttt{CSCOPE} : \mathcal{C}(S_2) : \texttt{PSCOPE} : \texttt{NOOP},\, 0 : e,\, \sigma,\, s,\, pc + l_b \rangle \\
              &\triangleright \langle \texttt{CSCOPE} : \mathcal{C}(S_2) : \texttt{PSCOPE} : \texttt{NOOP},\, e,\, \sigma,\, s,\, pc_{\texttt{CSCOPE}_2} \rangle
              && \text{(operational semantics)} \\
              &\triangleright \langle \mathcal{C}(S_2) : \texttt{PSCOPE} : \texttt{NOOP},\, e,\, \tilde{\sigma_c},\, \tilde{s_c},\, pc_{\texttt{CSCOPE}_2} + 1\rangle
              && \text{(operational semantics)} \\
              &\triangleright^\ast \langle \texttt{PSCOPE} : \texttt{NOOP},\, e',\, \underline{\sigma_c}' ,\, s'_c,\, pc_{\texttt{CSCOPE}_2} + l_{S_2} + 1 \rangle
              && \text{(induction hypothesis)} \\
              &\triangleright \langle \texttt{NOOP},\, e',\, \underline{\sigma}' ,\, s',\, pc_{\texttt{CSCOPE}_2} + l_{S_2} + 2 \rangle
              && \text{(operational semantics)} \\
              &\triangleright \langle \epsilon,\, e',\, \underline{\sigma}' ,\, s',\, pc_{\texttt{CSCOPE}_2} + l_{S_2} + 3 \rangle
              && \text{(operational semantics)} \\
              \end{aligned}
          \]
                
        \end{itemize}

    \item \textbf{case 9:} $P = (\texttt{while } b \texttt{ do } S \texttt{ od})$\\

        \begin{itemize}
            \item \textbf{subcase 9.1:} $\mathcal{B}(b, \sigma, s) = 0$

            Assume

            \[
                  \begin{aligned}
                  \langle P,\, \sigma,\, s \rangle
                  &=\langle \texttt{while } b \texttt{ do } S \texttt{ od},\, \sigma,\, s \rangle \\
                  &\rightarrow \langle \sigma ,\, s \rangle
                  && [\text{while}^0_{\text{os}}]   \\
                  \end{aligned}
              \]

            then,

            \[
              \begin{aligned}
              \langle \mathcal{C}(P),\, e,\, \sigma,\, s,\, pc \rangle
              &= \langle \mathcal{C}(\texttt{while } b \texttt{ do } S \texttt{ od}),\, e,\, \sigma,\, s,\, pc \rangle \\
              &= \langle \mathcal{CB}(b) : \texttt{JZ}\; pc_{\texttt{NOOP}} : \texttt{CSCOPE} : \mathcal{C}(S) : \texttt{PSCOPE} : \texttt{GOTO}\; pc_b : \texttt{NOOP},\, e,\, \sigma,\, s,\, pc \rangle
              && \text{(definition of $\mathcal{C}$)} \\   
              &\triangleright^\ast \langle \texttt{JZ}\; pc_{\texttt{NOOP}} : \texttt{CSCOPE} : \mathcal{C}(S) : \texttt{PSCOPE} : \texttt{GOTO}\; pc_b : \texttt{NOOP},\, \mathcal{B}(b, \sigma, s) : e,\, \sigma ,\, s,\, pc + l_b \rangle
              && \text{(correctness of $\mathcal{CB}$)} \\
              &= \langle \texttt{JZ}\; pc_{\texttt{NOOP}} : \texttt{CSCOPE} : \mathcal{C}(S) : \texttt{PSCOPE} : \texttt{GOTO}\; pc_b : \texttt{NOOP},\, 0 : e,\, \sigma ,\, s,\, pc + l_b \rangle \\
              &\triangleright \langle \texttt{NOOP},\, e,\, \sigma ,\, s,,\, pc_{\texttt{NOOP}} \rangle
              && \text{(operational semantics)} \\
              &\triangleright \langle \epsilon,\, e,\, \sigma ,\, s,,\, pc_{\texttt{NOOP}} + 1 \rangle
              && \text{(operational semantics)} \\
              \end{aligned}
          \]


          \item \textbf{subcase 9.2:} $\mathcal{B}(b, \sigma, s) = 1$

          \[
              \begin{aligned}
              \langle P,\, \sigma,\, s \rangle
              &=\langle \texttt{while } b \texttt{ do } S \texttt{ od},\, \sigma,\, s \rangle \\
              &\rightarrow \langle S\texttt{;endscope} \texttt{;while } b \texttt{ do } S \texttt{ od},\, \tilde{\sigma_c},\, \tilde{s_c} \rangle
              && [\text{while}^1_{\text{os}}]   \\
              &\rightarrow^\ast \langle \texttt{endscope} \texttt{;while } b \texttt{ do } S \texttt{ od},\, \sigma'_c,\, s'_c \rangle
              && ([\text{comp}^1_{\text{os}}] \,,\, [\text{comp}^2_{\text{os}}])   \\
              &\rightarrow \langle \texttt{while } b \texttt{ do } S \texttt{ od},\, \underline{\sigma}',\, s' \rangle
              && ([\text{scope}_{\text{os}}] \,,\, [\text{comp}^1_{\text{os}}] \,,\, [\text{comp}^2_{\text{os}}])   \\
              &&    \\
              \end{aligned}
          \]

          and assume for an induction hypothesis, that is, the compilation of $S$ is correct, then

          \[
              \begin{aligned}
              \langle \mathcal{C}(P),\, e,\, \sigma,\, s,\, pc \rangle
              &= \langle \mathcal{C}(\texttt{while } b \texttt{ do } S \texttt{ od}),\, e,\, \sigma,\, s,\, pc \rangle \\
              &= \langle \mathcal{CB}(b) : \texttt{JZ}\; pc_{\texttt{NOOP}} : \texttt{CSCOPE} : \mathcal{C}(S) : \texttt{PSCOPE} : \texttt{GOTO}\; pc_b : \texttt{NOOP},\, e,\, \sigma,\, s,\, pc \rangle
              && \text{(definition of $\mathcal{C}$)} \\   
              &\triangleright^\ast \langle \texttt{JZ}\; pc_{\texttt{NOOP}} : \texttt{CSCOPE} : \mathcal{C}(S) : \texttt{PSCOPE} : \texttt{GOTO}\; pc_b : \texttt{NOOP},\, \mathcal{B}(b, \sigma, s) : e,\, \sigma ,\, s,\, pc + l_b \rangle
              && \text{(correctness of $\mathcal{CB}$)} \\
              &= \langle \texttt{JZ}\; pc_{\texttt{NOOP}} : \texttt{CSCOPE} : \mathcal{C}(S) : \texttt{PSCOPE} : \texttt{GOTO}\; pc_b : \texttt{NOOP},\, 1 : e,\, \sigma ,\, s,\, pc + l_b \rangle \\
              &\triangleright \langle \texttt{CSCOPE} : \mathcal{C}(S) : \texttt{PSCOPE} : \texttt{GOTO}\; pc_b : \texttt{NOOP},\, 1 : e,\, \sigma ,\, s,\, pc + l_b + 1 \rangle
              && \text{(operational semantics)} \\
              &\triangleright \langle \mathcal{C}(S) : \texttt{PSCOPE} : \texttt{GOTO}\; pc_b : \texttt{NOOP},\, e,\, \tilde{\sigma_c},\, \tilde{s_c},\, pc + l_b + 2 \rangle
              && \text{(operational semantics)} \\
              &\triangleright^\ast \langle \texttt{PSCOPE} : \texttt{GOTO}\; pc_b : \texttt{NOOP},\, e',\, \sigma'_c,\, s'_c,\, pc + l_b + l_S + 2 \rangle
              && \text{(induction hypothesis)} \\
              &\triangleright \langle \texttt{GOTO}\; pc_b : \texttt{NOOP},\, e',\, \underline{\sigma}',\, s',\, pc + l_b + l_S + 3 \rangle
              && \text{(operational semantics)} \\
              &\triangleright \langle \mathcal{CB}(b) : \texttt{JZ}\; pc_{\texttt{NOOP}} : \texttt{CSCOPE} : \mathcal{C}(S) : \texttt{PSCOPE} : \texttt{GOTO}\; pc_b : \texttt{NOOP},\, e',\, \underline{\sigma}',\, s',\, pc_b \rangle
              && \text{(operational semantics)} \\
              &= \langle \mathcal{CB}(b) : \texttt{JZ}\; pc_{\texttt{NOOP}} : \texttt{CSCOPE} : \mathcal{C}(S) : \texttt{PSCOPE} : \texttt{GOTO}\; pc_b : \texttt{NOOP},\, e',\, \underline{\sigma}',\, s',\, pc \rangle \\
              &= \langle \mathcal{C}(\texttt{while } b \texttt{ do } S \texttt{ od}),\, e',\, \underline{\sigma}',\, s',\, pc \rangle \\
              \end{aligned}
          \]
        \end{itemize}

        However, the proofs above are not enough to establish the correctness of the compiled \texttt{while}-statements, because they only demonstrate one step of iteration. Nonetheless, with a simple induction, we can generalize this argument to claim that at any step of iteration, the compilation is correct, as the proofs above show that at any given program state and scope, the compilation of one iteration is correct. To put it more explicitly, the cases of 0 iterations (\textbf{subcase 6.1}) and 1 iteration (\textbf{subcase 6.2}) are proven to be correct. Assuming that $i$ iterations are compiled correctly, \textbf{subcase 6.2} also prove that $i+1$ iterations are compiled correctly. Ultimately, one can certainly say that any terminating \texttt{SelVeri} loop is compiled correctly.


  \item \textbf{case 10:} $P = \texttt{call}\; f(a_1 \dots, a_n)$

  Assume that $s'(x_i) = t_i$ and $\sigma'(x_i) = \mathcal{A}(a_i, \sigma, s)$ for $i = 1,2,\dots,n$

    \[
      \begin{aligned}
      \langle P,\, \sigma,\, s \rangle
      &=\langle \texttt{call}\; f(a_1 \dots, a_n),\, \sigma,\, s \rangle \\
      &\rightarrow \langle x_1 : t_1 \texttt{;} x_1 \,\texttt{:=}\, v_1\, \texttt{;} \dots \texttt{;} x_n : t_n \texttt{;} x_n \,\texttt{:=}\, v_n\, \texttt{;} S_f,\, \tilde{\sigma}', \tilde{s}' \rangle
      && [\text{call}_{\text{os}}]   \\
      &\rightarrow^\ast \langle S_f,\, \sigma' ,\, s' \rangle
      && (\text{operational semantics})   \\
      &\rightarrow^\ast \langle \texttt{return}\; a,\, \sigma'' \footnotemark ,\, s'' \rangle
      && (\text{operational semantics},\, \textit{return rule})   \\
      &\rightarrow \langle \sigma[\texttt{retvar} \mapsto \mathcal{A}(a, \sigma'', s'')],\; s[\texttt{retvar} \mapsto s''(a)] \rangle
      && [\text{ret}_{\text{os}}] \\
      \end{aligned}
  \]

by the induction hypothesis, the compilation of $S_f$ is correct, since the execution of the function body needs lesser steps than the whole function call; then

\[
  \begin{aligned}
  \langle \mathcal{C}(P),\, e,\, \sigma,\, s,\, pc,\, ras \rangle
  &= \langle \mathcal{C}(\texttt{call}\; f(a_1 \dots, a_n)),\, e,\, \sigma,\, s,\, pc \rangle \\
  &= \langle \mathcal{CA}(a_n) : \dots : \mathcal{CA}(a_1) : \texttt{CALL}\; f,\, e,\, \sigma,\, s,\, pc,\, ras \rangle
  && (\text{definition of $\mathcal{C}$})   \\
  &\triangleright^\ast \langle  \texttt{CALL}\; f,\, \mathcal{A}(a_1, \sigma, s) : \dots : \mathcal{A}(a_n, \sigma, s) : e,\, \sigma,\, s,\, pc+n,\, ras \rangle
  && (\text{correctness of $\mathcal{CA}$})   \\
  &\triangleright \langle  P[pc_f],\, \mathcal{A}(a_1, \sigma, s) : \dots : \mathcal{A}(a_n, \sigma, s) : e,\, \sigma,\, s,\, pc_f,\, (pc+n+1) : ras \rangle
  && (\text{operational semantics})   \\
  &= \langle  \texttt{FUNCENV}\; f, t_r : \texttt{DECL}\;t_1\,x_1 : \texttt{STORE}\; x_1 :  \dots : \texttt{DECL}\;t_1\,x_1 : \texttt{STORE}\; x_n : \mathcal{C}(S_f),\, \mathcal{A}(a_1, \sigma, s) : \dots : \mathcal{A}(a_n, \sigma, s) : e,\, \sigma,\, s,\, pc_f,\, (pc+n+1) : ras \rangle
  && (\text{definition of $\mathcal{C}$})   \\
  &\triangleright \langle  \texttt{DECL}\;t_1\,x_1 : \texttt{STORE}\; x_1 :  \dots : \texttt{DECL}\;t_1\,x_1 : \texttt{STORE}\; x_n : \mathcal{C}(S_f),\, \mathcal{A}(a_1, \sigma, s) : \dots : \mathcal{A}(a_n, \sigma, s) : e,\, \tilde{\sigma}', \tilde{s}',\, pc_f + 1,\, (pc+n+1) : ras \rangle
  && (\text{operational semantics})   \\
  &\triangleright^\ast \langle \mathcal{C}(S_f),\, e,\, \sigma',\, s',\, pc_f + 2n+1,\, (pc+n+1) : ras \rangle
  && (\text{operational semantics})   \\
  &\triangleright^\ast \langle \mathcal{C}(\texttt{return}\; a),\, e,\, \sigma'',\, s'',\, pc_f + l + 2n+1,\, (pc+n+1) : ras \rangle
  && (\text{induction hypothesis},\, \textit{return rule})   \\
  &=\langle \mathcal{CA}(a) : \texttt{RET},\, e,\, \sigma'',\, s'',\, pc_f + l + 2n+1,\, (pc+n+1) : ras \rangle
  && (\text{definition of $\mathcal{C}$})   \\
  &\triangleright^\ast \langle \texttt{RET},\, \mathcal{A}(a, \sigma'', s'') : e,\, \sigma'',\, s'',\, pc_f + l' + 2n+1,\, (pc+n+1) : ras \rangle
  && (\text{correctness of $\mathcal{CA}$})   \\
  &\triangleright \langle P[pc+n+1:],\, e,\, \sigma[\texttt{retvar} \mapsto \mathcal{A}(a, \sigma'', s'')],\; s[\texttt{retvar} \mapsto s''(a)],\, pc+n+1,\, ras \rangle
  && (\text{operational semantics})   \\
  
  \end{aligned}
\]

Similar to our proof of the correctness of \texttt{while}-statements, this assumes a finite depth of recursion (0, if there is no recursion) as we have assumed termination. \\

\footnotetext{Note that we do not cover the case where the execution of $S_f$ leads to an erroneous state, as the last semantic rule in the section \textbf{2.2.5} of this document and [$\text{err}_{\text{os}}$] rule of high-level semantics ensure the correctness of the translation in this case.}

  \item \textbf{case 11:} $P = read$

  \item \textbf{case 12:} $P = S_1\texttt{;}S_2$ for some statements $S_1$, $S_2$\\

    Assume

    \[
      \begin{aligned}
      \langle P,\, \sigma,\, s \rangle
      &=\langle S_1\texttt{;}S_2,\, \sigma,\, s \rangle \\
      &\rightarrow^\ast \langle S_2,\, \underline{\sigma}' ,\, s' \rangle
      && ([\text{comp}^1_{\text{os}}] \,,\, [\text{comp}^2_{\text{os}}])   \\
      &\rightarrow^\ast \langle \underline{\sigma}'' ,\, s'' \rangle
      && ([\text{comp}^1_{\text{os}}] \,,\, [\text{comp}^2_{\text{os}}]) \\
      \end{aligned}
  \]
  
  and, assume for an induction, that is the compilation of $S_1$ and $S_2$ are correct, then

  \[
      \begin{aligned}
      \langle \mathcal{C}(P),\, e,\, \sigma,\, s,\, pc \rangle
      &= \langle \mathcal{C}(S_1\texttt{;}S_2),\, e,\, \sigma,\, s,\, pc \rangle \\
      &= \langle \mathcal{C}(S_1) : \mathcal{C}(S_2) ,\, e,\, \sigma,\, s,\, pc \rangle
      && \text{(definition of $\mathcal{C}$)} \\   
      &= \langle \mathcal{C}(S_2) ,\, e',\, \underline{\sigma}' ,\, s',\, pc + l_{S_1} \rangle
      && \text{(induction hypothesis 1)} \\   
      &= \langle \epsilon ,\, e'',\, \underline{\sigma}'' ,\, s'',\, pc + l_{S_1} + l_{S_2} \rangle
      && \text{(induction hypothesis 2)} \\ 
      \end{aligned}
  \]

  
    

\end{itemize}

Observe that at each case, the resulting program state and scope tuple is the equal for \texttt{SelVeri} and \texttt{SelVerIR}, the difference between initial and final program counters are equal to $l_p$ and $ras$ is not altered, therefore concluding \footnote{Note that we need not prove the correctness of \texttt{return}-statements separately; as they are always a part of function calls, case 7 also establishes the correctness of \texttt{return}-statements} the proof.
\\
\end{proof}

\newpage
Cautious readers may notice that some (sub)cases including erroneous program states (either as input or output) is omitted in our correctness proofs. One of the following reasons is the justification for this omittance:

\begin{enumerate}
    \item Such a case cannot be encountered in the range of neither $\mathcal{CA}$ nor $\mathcal{CB}$ nor $\mathcal{C}$.

    \item Such an erroneous case can be detected during compile-time, therefore no \texttt{SelVerIR} code inheriting the error can be generated.

    \item After a finite number of non-erroneous execution steps, which are proven correct in this section, program reaches an error state; but by section \textbf{2.2.5} of this document and the operational semantic rule of high-level \texttt{SelVeri} language $[ \text{err}_{\text{os}} ]$ the remaining finite execution is also correct.
\end{enumerate}

No matter the reason, as no statically-correct \texttt{SelVeri} program includes any of the omitted cases, the correctness proofs above are sufficient for establishing the correctness of our implementation. \\

\vspace{1em}
\section{Divergent programs}

If a \texttt{SelVeri} program $P$ terminates after a finite number of steps, it is called \textit{convergent}; otherwise, a non-terminating $P$ is called \textit{divergent}. \texttt{SelVeri} has two possible conditions for non-termination:

\begin{enumerate}
    \item Infinite \texttt{while}-loops
    \item Infinite recursion
\end{enumerate}

Therefore any infinite program, after the execution of a finite subprogram, is an infinite repetition of a (possibly compositional) statement, which is either a loop body, function body or a composition of function bodies; whose execution takes a finite number of steps. We call this fact as the \textit{decomposition lemma}, which will play a very important role for our correctness proof for divergent programs. Nonetheless, our formalization so far is not capable of speaking about infinite execution traces, therefore an extension is ultimately necessary \footnote{Note that we will not be considering erroneous programs in this section, as they are finite by their nature.}. We will be utilizing ordinal numbers for this purpose.\\

\subsection{Ordinal program states}

Let $\alpha$ be an (possibly infinite) ordinal number; then, given any configuration, we have:

\[
\langle P,\, \sigma,\, s \rangle \rightarrow^\alpha \langle P',\, \sigma^P_\alpha,\, s^P_\alpha \rangle
\] where $\rightarrow^\alpha$ means execution of the program $\alpha$-many steps. \\

If $\alpha$ is a finite ordinal (i.e. a natural number), then we have already formalized what the transition above means. In the infinite ordinal case, $\sigma^P_\alpha$ is defined as follows:

\[
\sigma^P_\alpha(x, s^P_\alpha)  = 
\begin{cases}
    v & \text{if $v \in \textbf{Val} \;\land\; \exists \beta < \alpha .\; \forall \gamma .\; \beta < \gamma < \alpha \implies \sigma^P_\gamma(x, s^P_\gamma) = v$} \\
    \mu & \text{otherwise} \\
\end{cases}
\] where $\mu$ is a special constant, denoting the value \textit{unknown}. A reason for having the value unknown may be oscillation or indefinite growth/decrease. If the value of any variable $x$ is \textit{known} after executing $\alpha$-many steps, then we say that $x$, \textit{$\alpha$-stabilizes}. \\

Moreover, since \texttt{SelVeri} allows a variable to be declared only once in a non-erroneous program, we need not extend the formalization of scopes. Throughout this section, we will restrict our attention on how the program states behave under infinite traces.

\subsection{Transfinite correctness of divergent programs}

Now, as our formalization enhanced, we can attempt to prove the correctness of divergent \texttt{SelVeri} programs. More formally, we need to show that; given any configuration and an ordinal $\alpha$ 
\[
\langle P,\, \sigma,\, s \rangle \; \rightarrow^\alpha \; \langle P',\, \sigma^P_\alpha,\, s^P_\alpha \rangle \quad \implies \quad 
\langle \mathcal{C}(P),\, e,\, \sigma,\, s,\, pc,\, ras \rangle \; \triangleright^\beta \; \langle \mathcal{C}(P'),\, e',\, \sigma^P_\alpha,\, s^P_\alpha,\, pc',\, ras' \rangle 
\] where $\beta = k  \cdot \alpha$ \footnote{Because the compilation function $\mathcal{C}$ emits a finite number of commands per high-level statement.} for some constant non-zero $k \in \mathbb{N}$. 

\begin{proof} We will prove this by a transfinite induction on $\alpha$. But before starting the proof; observe that the small-step operational semantics of \texttt{SelVeri} and \texttt{SelVerIR} are defined exclusively via successor transition relations, meaning that any execution trace is a sequence strictly indexed by the natural numbers, $\mathbb{N}$. Consequently, the execution trace of any divergent program has an exact ordinal length of $\omega$, as the semantics do not define limit-ordinal transitions. Hence, we limit our attention to the ordinals that are lesser or equal to $\omega$ \footnote{We are not interested in hypercomputation here.}. We call this fact as \textit{$\omega$-bound of computation}.

    \begin{itemize}
        \item \textbf{The Base Case:} $\alpha = 0$. \\
        Zero steps taken, the states are trivially identical.

        \item \textbf{The Successor Case:} Assume that theorem holds for some ordinal $\gamma$ and let $\alpha = \gamma^+$. \\
        \textbf{Section 5} exactly proves the successor case. Note that, the proof in the \textbf{section 5} trivially implies the prefix correctness via the semantics of compositional statements.

        \item \textbf{The Limit Case:} Let $\alpha$ be a limit ordinal and assume that the theorem holds for all ordinals $\gamma < \alpha$. \\ By the \textit{$\omega$-bound of computation}, we have $\alpha = \beta = \omega$. This also means our induction hypothesis is equivalent to the correctness of convergent programs, which is already proven. Additionally, because the only reasons of divergence in a \texttt{SelVeri} program are infinite loops or recursions, after a finite number of steps, say $M$; by the decomposition lemma there must be an indefinitely executed statement $S_r$  whose longest possible execution takes a finite number of steps, say $m$. Given any variable $x$, consider two cases:

        \begin{enumerate}
            \item $x$ $\omega$-stabilizes. \\\\
            Then by the definition of stabilization, $\sigma^P_\omega(x, s^P_\omega) = v$ for some $v \in \textbf{Val}$, meaning that there exists a natural number $N$ such that for all natural numbers $n > N$, we have $\sigma^P_n(x, s^P_n) = v$. Now, let $K = \textit{max}(N, M)$ and observe that for all $k$ such that $K < k < \omega$, we have $\sigma^P_k(x, s^Pk) = v$ since $x$ is $\omega$-stabilizing. Also, by our choice, the $k^\text{th}$ step exactly lies in the indefinite repetition of $S_r$. This means that, after the step $K$; $S_r$ (or any sub-statement of $S_r$) either never assigns any value to $x$ or if it assigns, the value is strictly equal to $v$; otherwise the stabilization of $x$ would be violated. In other words, we have reached to the fixed-point \footnote{We are slightly drifting into the denotational semantics region.} of the value of $x$, under the semantics of $S_r$. \\
            Acknowledging the fact that $S_r$ by itself is a finite program, $\mathcal{C}(S_r)$ does not modify the value of $x$ after some step $K'$ which is greater than $K$ but still finite, as any compiled statement executes more (but still finite) steps than its high-level counterpart, by the correctness of finite compilation. This means that $x$ is also an $\omega$-stabilizing variable under the execution of $\mathcal{C}(P)$.

            \vspace{2em}
            
            \item $x$ does not $\omega$-stabilize. ($\sigma^P_\omega(x, s^P_\omega) = \mu$) \\\\
            Then, for an arbitrary finite step $N > M$, there exists some future step $n > N$ such that $\sigma^P_N(x, s^P_N) \neq \sigma^P_n(x, s^P_n)$. Let $i = \lceil (n- N)/m \rceil$ and let $S^N_r$ be the composition of $S_r$ by itself $i$-many times. (Similar to $S_r$, $S^N_r$ is also executed indefinitely.) Then by construction, both the $N^\text{th}$ and $n^\text{th}$ steps of execution exactly lie in some indefinite repetition of $S^N_r$, meaning that $S^N_r$ does not preserve the value of the variable $x$. By its definition, one execution of $S^N_r$ takes at most $i\cdot m$ steps, which is finite. By correctness of finite compilations (i.e. the induction hypothesis), $\mathcal{C}(S^N_r)$ also does not preserve the value of $x$. \\
            Observe that with the exact same manner, for any sufficiently large $N'$, we can    construct a statement $S^{N'}_r$ (therefore a $\mathcal{C}(S^{N'}_r)$ in the intermediate representation)  that does not preserve the value of $x$, hence $x$ also does not $\omega$-stabilize under the execution of $\mathcal{C}(P)$, i.e. the value of $x$ is $\mu$ in the IR counterpart as well.      
        \end{enumerate}

        Both of these cases ensure that the resulting program states after executing $P$ and $\mathcal{C}(P)$ for $\omega$-steps, are exactly equal; establishing the correctness of the compilation of divergent \texttt{SelVeri} programs.
        
    \end{itemize}
\end{proof}

\section{Turing-completeness of \texttt{SelVerIR}}

As \texttt{SelVeri} is a Turing-complete programming language that can be embedded into \texttt{SelVerIR} in a semantically-equivalent way, via the compilation function $\mathcal{C}$; \texttt{SelVerIR} is able to compute any computable function, hence is Turing-complete.

\end{document}
